%\documentclass{article}
%\documentclass[a4paper]{amsart}


\documentclass[11pt,a4paper]{article}%\newcommand{\correctspacing}  {\onehalfspacing} 

%\usepackage{showkeys}
%\usepackage{showkeys}

%\usepackage[notref,notcite]{showkeys}
\usepackage{amsmath,amsthm,amsfonts,amssymb,epsfig}
\usepackage{setspace}
\usepackage{xcolor}

\topmargin -0.0cm 
 \textheight 22.0cm

\oddsidemargin -0.cm \evensidemargin -0.cm
\textwidth 16cm

%{\singlespacing}

%\usepackage{a4}
\usepackage{amsmath, amsthm, mathscinet}
%\usepackage{amsthm}
\usepackage{amssymb}
\usepackage{
	comment,
	%thmtools, % enables listoftheorem and defining theorem
	enumitem, % enables a greater range of eumerate options
	%pgfplots, % plot generation package
	float, % images option H
	wrapfig,
	subcaption }
\usepackage{microtype}
%\usepackage{todonotes}
%\usepackage[unicode=true,bookmarks=true,bookmarksnumbered=true, breaklinks=false, pdfborder={0 0 0}, pdfborderstyle={}, colorlinks=true, urlcolor= blue, citecolor= blue, pdftitle={nested}]{hyperref}
\usepackage[applemac]{inputenc}
\usepackage{graphicx}
\usepackage{color}
\usepackage{tikz} \usetikzlibrary{trees}
%\usepackage{enumerate}
\usepackage{enumitem}
\setlist[enumerate,1]{label=(\roman*)}
%\usepackage{a4wide}

\usetikzlibrary{matrix}
%\topmargin -1.0cm
%\oddsidemargin 0in
%\evensidemargin 0in
%\textwidth 6.3 truein
%\textheight 9.1 truein

\makeatletter
\newcommand*\bigcdot{\mathpalette\bigcdot@{.6}}
\newcommand*\bigcdot@[2]{\mathbin{\vcenter{\hbox{\scalebox{#2}{$\m@th#1\bullet$}}}}}
\makeatother
\def\comment#1{\marginpar{\raggedright\scriptsize{\textcolor{purple}{#1}}}}
%\def\comment#1{}

%%%%%%%%%%%%%%%%%%%%%%%%%%%%%%%%%%%%%%%%%%%%%%%%%%%%%%%%%%%%%%%%%%%%%%%%%%%%%%%%%%%%%%%%%%%%%%%%%%%%%%%%%%%%%%%%%%%%%%%%%%%%%%%%%%%%%%%%%%%%%%%%%%%%%%%%%%%%%%%%%%%%%%%%%%%%%%%%%%%%%%%%%%%%%%%%%%%%%%%%%%%%%%%%%%%%%%%%%%%%%%%%%%%%%%%%%%%%%%%%%%%%%%%%%%%%
\newcommand{\cblue}{\color{red}}
\newcommand{\cred}{\color{blue}}
\newcommand{\cgreen}{\color{green}}
\newcommand{\cpurple}{\color{purple}}
\newcommand{\ccyan}{\color{cyan}}


\def\diam{\mathop{\rm diam}}
\def\argmin{\mathop{\rm arg\, min}}
\def\A{{\cal A}}
\def\B{{\cal B}}
\def\C{{\cal C}}
\def\D{{\cal D}}
\def\E{{\mathbb E}}
\def\F{{\mathcal F}}
\def\bbF{{\mathbb F}}
\def\bbW{{\mathbb W}}
\def\G{{\cal G}}
\def\H{{\cal H}}
\def\I{{\cal I}}
\def\J{{\cal J}}
\def\K{{\cal K}}
\def\L{{\mathcal L}}
\def\M{{\cal M}}
\def\N{{\cal N}}
\def\NN{\mathbb{ N}}
\def\O{{\cal O}}
\def\P{{\mathbb P}}
\def\Q{{\mathbb Q}}
\def\R{{\mathbb R}}
\def\S{{\cal S}}
\def\T{{\cal T}}
\def\U{{\cal U}}
\def\V{{\mathbb V}}
\def\W{{\cal W}}
\def\X{{\mathcal X}}
\def\x{{\cal x}}
\def\y{{\cal y}}
\def\Y{{\mathcal Y}}
\def\Z{{\cal Z}}
\def\QQ{{\cal Q}}
\def\PROB{{\mathbb P}}
\def\EXP{{\mathbb E}}
\def\Var{{\bf Var}}
\def\I\nd{{\mathbb I}}
\def\e{{\epsilon}}
%\def\l{{\ell}}
\def\bbT{{\mathbb T}}
\def\ci{\perp\!\!\!\perp}
\renewcommand{\subset}{\subseteq}
\renewcommand{\supset}{\supseteq}

\newtheorem{theorem}{Theorem}[section]
\newtheorem{acknowledgement}[theorem]{Acknowledgement}
\newtheorem{algorithm}[theorem]{Algorithm}
\newtheorem{axiom}[theorem]{Axiom}
\newtheorem{case}[theorem]{Case}
\newtheorem{claim}[theorem]{Claim}
\newtheorem{conclusion}[theorem]{Conclusion}
\newtheorem{conjecture}[theorem]{Conjecture}
\newtheorem{corollary}[theorem]{Corollary}
\newtheorem{criterion}[theorem]{Criterion}
\newtheorem{exercise}[theorem]{Exercise}
\newtheorem{lemma}[theorem]{Lemma}
\newtheorem{notation}[theorem]{Notation}
\newtheorem{problem}[theorem]{Problem}
\newtheorem{proposition}[theorem]{Proposition}
\newtheorem{solution}[theorem]{Solution}
\newtheorem{definition}[theorem]{Definition}
\newtheorem{condition}[theorem]{Condition}
\newtheorem{assumption}[theorem]{Assumption}
\newtheorem{remark}[theorem]{Remark}
\newtheorem{example}[theorem]{Example}
\newtheorem{summary}[theorem]{Summary}

%\numberwithin{equation}{section}

\renewcommand{\epsilon}{\varepsilon}
\newcommand{\nd}{\text{nd}}

%\renewcommand{\langle}{[}
%\renewcommand{\rangle}{]}

\newcommand{\AW}{\mathcal{AW}}
\newcommand{\eps}{\varepsilon}

\newcommand{\id}{\text{Id}}

\newcommand{\UW}{\mathcal{AW}}


\newcommand{\OB}{Ob{\l}{\'o}j}%{Ob{\l}{\'o}j} 
\newcommand{\HL}{Henry-Labord{\`e}re}

\newcommand{\pihc}{the PI and his collaborators}
\newcommand{\wat}{weak adapted topology}
\numberwithin{equation}{section}

\newcommand{\Couplings}{\text{Cpl}}

\newcommand{\cpl}{\Couplings}
\newcommand{\cpla}{\Couplings_c}
\newcommand{\cplba}{\Couplings_{bc}}
%\newcommand{\F}{\mathscr F}
\newcommand{\WA}{\mathcal W}
\newcommand{\AWA}{\mathcal {AW}}
\newcommand{\CWA}{\mathcal {CW}}
\newcommand{\SWA}{\mathcal {SCW}}
\newcommand{\SCWA}{\mathcal{SCW}}
\newcommand{\IWA}{\mathcal{IW}}
\newcommand{\ND}{\mathcal{ND}}

\newcommand{\Wp}{\mathcal W_p}
\newcommand{\AWp}{\mathcal A \mathcal W_p}
\newcommand{\CWp}{\mathcal C \mathcal W_p}
\newcommand{\SCWp}{\mathcal S \mathcal C \mathcal W_p}

\newcommand{\WAT}{weak adapted topology}

\title{From Bachelier to Dupire via Optimal Transport
}

%[Backhoff, Bartl, Beiglb\''ock, Eder]
 \author{Mathias Beiglb\"ock, Gudmund Pammer,
Walter Schachermayer}
\begin{document}


\maketitle
\begin{abstract} 
Famously mathematical finance was started by Bachelier  in his 1900 PhD thesis where -- among many other achievements -- he also provides a formal derivation of the Kolmogorov forward equation. This forms also the basis for Dupire's (again formal) solution to the problem of finding an arbitrage free model calibrated to the volatility surface. 
 The later result has  rigorous counterparts in the theorems of Kellerer and Lowther. In this survey article we revisit these hallmarks of stochastic finance,  highlighting the role played by some optimal transport results in this context.

%{\cpurple MB: Der abstract ist bis jetzt auch nur ein Denkanstoss...}

{\em Keywords:}  
%\begin{keywords}
Bachelier, Dupire's formula, Kellerer's theorem,  optimal transport, martingales, peacocks.
%\end{keywords}
\end{abstract} 




%{\bf AMS subject classifications.}  



%\section{Introduction}

\section{Bachelier's work relating Brownian motion to mass transport and the heat equation} \label{ahnbwrbmtthe}



%%%%%%%%%%%%%%%%%%%%%%%%%%%%%%%%%%%%%%%%%%%%%%%%%%%%%%%%%%%%%%%%%%%%%%%%%%%%%%%%%%%%%%%%%%%%%%%%%%%%%%%%%%%%%%%%%%%%%%%%%%%%%%%%%%%%%%%%%%%%%%%%%%%%%%%%%%%%%%%%%%%%%%%%%%%%%%%%%%%%%%%%%%%%%%%%%%%%%%%%%%%%%%%%%%%%%%%%%%%&&&&&&&&&&&&&&&&&&&&&&&&&&&&&&&&&

In this section, which is mainly dedicated to the historic point of view,  {we follow \cite{Sc03}} and point out that Bachelier already had some thoughts on \textit{``horizontal transport of probability measures''} in his dissertation \textit{``Th{\'e}orie de la sp{\'e}culation''} \cite{Ba00}, which he defended in Paris in 1900.


In this thesis he was the first to consider a \textit{mathematical model} of Brownian motion. Bachelier argued using infinitesimals by visualizing Brownian motion $(W(t))_{t \geqslant 0}$ as an infinitesimal version of a random walk. His 19th century style argument runs as follows. Suppose that the grid in space is given by 
\begin{equation} \label{grid}
\ldots, \ x_{n-2}, \ x_{n-1}, \ x_{n}, \ x_{n+1}, \ x_{n+2}, \ \ldots
\end{equation}
having the same (infinitesimal) distance $\Delta x = x_{n}-x_{n-1}$, for all $n$, and such that at time $t$ these points have (infinitesimal) probabilities 
\begin{equation} \label{gridp}
\ldots, \ p_{n-2}^{t}, \ p_{n-1}^{t}, \ p_{n}^{t}, \ p_{n+1}^{t}, \ p_{n+2}^{t}, \ \ldots
\end{equation}
under the random walk under consideration. What are the probabilities 
\begin{equation} \label{gridpp}
\ldots, \ p_{n-2}^{t+\Delta t}, \ p_{n-1}^{t+\Delta t}, \ p_{n}^{t+\Delta t}, \ p_{n+1}^{t+\Delta t}, \ p_{n+2}^{t+\Delta t}, \ \ldots
\end{equation}
of these points at time $t+\Delta t$? 

The random walk moves half of the mass $p_{n}^{t}$, sitting on $x_{n}$ at time $t$, to the point $x_{n+1}$. En revanche, it moves half of the mass $p_{n+1}^{t}$, sitting on $x_{n+1}$ at time $t$, to the point $x_{n}$. We thus may calculate the net difference between $p_{n}^{t}/2$ and $p_{n+1}^{t}/2$, which Bachelier identifies with 
\begin{equation}
-\tfrac{1}{2} \, \frac{\partial p^{t}}{\partial x}(x), 
\end{equation}
where we let $x = x_n = x_{n+1}$ which is legitimate for Bachelier as $x_n$ and $x_{n+1}$ only differ by an infinitesimal.

This amount of mass is transported from the interval $(-\infty,x_{n}]$ to $[x_{n+1},\infty)$ during the time interval $ (t,t+\Delta t)$. In Bachelier's own words, this is very nicely captured by the following quote from his thesis: 

\medskip

\textit{``Each price $x$ during an element of time radiates towards its neighboring price an amount of probability proportional to the difference of their probabilities. I say proportional because it is necessary to account for the relation of $\Delta x$ to $\Delta t$. The above law can, by analogy with certain physical theories, be called the law of radiation or diffusion of probability.''}

\medskip

Passing formally to the continuous limit and -- using today's terminology -- denoting by 
\begin{equation} \label{dist}
P_t(x) = \int_{-\infty}^{x} p_t(z) \, \textnormal{d}z
\end{equation}
the distribution function associated to the Gaussian density function $p_t(x)$, Bachelier thus deduces in this intuitively convincing way the relation
\begin{equation} \label{lbhef}
\frac{\partial P}{\partial t} = \frac{1}{2} \frac{\partial p}{\partial x},
\end{equation}
where we have normalized the relation between $\Delta x$ and $\Delta t$ to obtain the constant $1/2$. By differentiating 
(\ref{lbhef}) with respect to $x$ one obtains the usual heat equation 
\begin{equation} \label{heflbhef}
\frac{\partial p}{\partial t} = \frac{1}{2} \frac{\partial^{2} p}{\partial x^{2}}
\end{equation}
for the density function $p(t,x)$. 
Of course, the heat equation was known to Bachelier, and he notes regarding (\ref{heflbhef})
 \textit{``C'est une {\'e}quation de Fourier.''}

%\begin{remark}\label{ForwardRemark}
Bachelier has thus derived, on a formal level, the Kolmogorov forward equation, also known as Fokker-Planck-equation, for the propagation of a probability density $p$ under Brownian motion. The forward equation will also play an important role subsequently and we take the opportunity to note that Bachelier's argument can equally well be applied to the more general process with increments $dX_t = \sigma(t,X_t)\, dW_t$ to arrive at the PDE
\begin{equation}\label{Fokker0}
\frac{\partial}{\partial t} P= \frac{1}{2} \frac{\partial}{\partial x} \left(\sigma^2 p\right),
\qquad \qquad
\frac{\partial}{\partial t}p = \frac12 \frac{\partial^2}{\partial x^2}\left(\sigma^2 p\right).
\end{equation}
%\end{remark}







But let us still remain with the form (\ref{lbhef})
 of the heat equation and analyze its message in terms of \textit{``horizontal transport of probability measures''}. One may ask: what is the ``velocity field'', acting on the set of probabilities on $\R$, which moves the probability density $p_t(\cdot)$ to the probability density $p_{t+dt}(\cdot)$? Following Bachelier's intuition and keeping in mind that the mass sitting at time $t$ on $x$ equals $p_t(x)$, the velocity of this move at the point $x$ has to be equal to
 
 \begin{equation} \label{score}
 -\frac{1}{2} \frac{\frac{\partial p_{t}}{\partial x}(x)}{p_t(x)} 
\end{equation}
which has the natural interpretation as the ``speed'' of the horizontal  transport induced by $p_t(x)$. We thus encounter \textit{in nuce} the  ``score function'' $\frac{p'_t(x)}{p_t(x)} = \frac{\nabla p_t(x)}{p_t(x)}$ where the nabla notation $\nabla$ indicates that this is a vector field which makes perfect sense in the $n$-dimensional case too.

At this stage we can relate Bachelier's work with the more recent notion of the \textit{Wasserstein metric} $\mathcal{W}_2 (\cdot,\cdot)$, at least intuitively and at an infinitesimal level. One may ask: what is the necessary kinetic energy needed to transport $p_t(\cdot)$ to $p_{t+dt}(\cdot)$? Knowing the speed (\ref{score}) and the usual formula for the kinetic energy, we obtain for the Wasserstein distance between the two infinitesimally close probabilities $p_t$ and $p_{t+dt}$

\begin{equation}\label{Fisher}
 \frac{\mathcal{W}_2 \left(p_t,p_{t+dt}\right)}{dt} := \frac{1}{2} \left(\int_{\R} \left(\frac{p'_t(x)}{p_t(x)}\right)^2 dx \right)^{\frac{1}{2}}
\end{equation}

For a formal definition of the Wasserstein distance $\mathcal{W}_2 (\cdot,\cdot)$ we refer, e.g., to \cite{Vi09}. While for the finite version of the Wasserstein distance between two probability measures one has to find an optimal transport plan, the situation is simpler -- and very pleasant -- in the case of the infinitesimal transport induced by the vector field (\ref{score}). This infinitesimal transport is automatically optimal.
Intuitively this corresponds to the geometric insight in the one-dimensional case that the transport lines of infinitesimal length cannot cross each other. For a thorough treatment of the geometry of absolutely continuous curves of probabilities such as $\left(p_t(\cdot)\right)_{t \ge 0}$ above we refer to the lecture notes \cite{AmGiSa08}. 





\medskip

%Let us go one step beyond Bachelier's thoughts and consider the relation of the above infinitesimal Wasserstein transport to time reversal (which Bachelier had not yet considered in his solitary exploration of Brownian motion). Visualizing again the grids (\ref{gridp}) and (\ref{gridpp}),
%a moment's reflection reveals that the transport from $p^{t+\Delta t}$ to $p^{t}$, i.e., in reverse direction, is accomplished by going from $x_{n}$ to $x_{n+1}$ with probability $\frac{1}{2} + \frac{p'(t,x)}{p(t,x)} \, \textnormal{d}x$, and from $x_{n+1}$ to $x_{n}$ with probability $\frac{1}{2} - \frac{p'(t,x)}{p(t,x)} \, \textnormal{d}x$, with the identifications $x = x_{n} = x_{n+1}$, and $\textnormal{d}x = \Delta x$. In other words, the above Brownian motion $(W(t))_{t \geqslant 0}$, considered in reverse direction, $(W(T-s))_{0 \leqslant s \leqslant T}$ is \textit{not} a Brownian motion any more, as the transition probabilities do not equal $(1/2,1/2)$. Rather, one has to correct these probabilities by a term which --- once again --- involves our familiar score function $\nabla p(t,x) / p(t,x)$ from (\ref{score})
%above. At this stage, it should come as no surprise, that the passage to reverse time is closely related to the Wasserstein transport induced by $p(t,x)$. This observation paves the way to an interplay of time reversal of diffusion processes with the notions of Wasserstein distance, Fisher information, and decay of entropy. We refer to \cite{KaScTs20} for more on this theme as well as an account of the development of these ideas since the work by Schr\''odinger and Nelson.

\medskip

We finish the section by returning to Bachelier's thesis. The \textit{rapporteur} of Bachelier's dissertation was no lesser a figure than Henri Poincar{\'e}. Apparently he was aware of the enormous potential of the section \textit{``Rayonnement de la probabilit{\'e}''} in Bachelier's thesis, when he added to his very positive report the handwritten phrase: \textit{``On peut regretter que M. Bachelier n'ait pas d{\'e}velopp{\'e} davantage cette partie de sa th{\`e}se.''} That is: One might regret that Mr. Bachelier did not develop further this part of his thesis. Truly prophetic words!




\section{Dupire's Formula} \label{Dupire}

We now turn to a well-known and more recent topic in Mathematical Finance continuing the early achievements of Bachelier.

Suppose that in a financial market we know the prices of ``many'' European options on a given (highly liquid) stock $S$. What can we deduce from this data about the prices of exotic, i.e., path-dependent options?

This question leads to the following mathematical idealization: suppose we know the prices of \textit{all} European call options, i.e., the price $C(t,x)$ of every call option with strike price $x$ and maturity $t$, for every $0 \le t \le T$ and $x \in \R_+$. Our task is to analyze the set of all possible (local) martingale measures for the stock price processes which are compatible with this data.
Once we have a hand on the relevant set of martingale measures, we can price arbitrary exotic options by taking expectations.

To make the question meaningful, it is a good idea to restrict the class of processes under consideration, e.g., to continuous, Markovian martingales.

We make the economically meaningful assumptions that the function $(t,x) \mapsto C(t,x)$ is sufficiently smooth, as well as strictly convex in the variable $x$ and strictly increasing in the variable $t$, to allow for the subsequent formal manipulations. 

%Later we shall have a closer look at precise mathematical statements.

The first observation is that the knowledge of $C(t,x)$, for $0\le t\le T$ and $x>0$ is tantamount to the knowledge of the marginal probabilities $(\mu_t)_{0 \le t \le T}$ of the underlying stock price process under a martingale measure which determines the prices, via the formula 
\begin{equation} \label{eq3.1}
C(t,x) = \E_{\mu_t}[(S_t-x)_+].
\end{equation}
This observation goes back to the Breeden and Litzenberger \cite{BrLi78}.

If the measures $\mu_t$ are absolutely continuous with respect to Lebesgue measure with a continuous density function $p_t(x)$ then (\ref {eq3.1}) amounts to the relation
\begin{equation}\label{eq3.2}
p_t(x) = C_{xx}(t,x),\qquad x>0,
\end{equation}
as one verifies via integration by parts.

In a very influential and highly cited paper from 1994 \cite{Du94} (compare also the work of Derman and Kani \cite{DeKa94}) Dupire  considered diffusion processes of the form 
\begin{equation}\label{eq3.3}
\frac{dS_t}{S_t}=\sigma(t,S_t)\, dW_t,\qquad 0\le t \le T,
\end{equation}
where the ``local volatility'' $\sigma(\cdot,\cdot)$ is modeled as a {\it deterministic} function  of $t$ and $x$, and $W_t$ is a Brownian motion adapted to its natural filtration $(\mathcal F_t)_{0 \le t \le T }$. It turns out that there is the beautiful and strikingly simple ``Dupire's formula'' which relates $\sigma(\cdot,\cdot)$ to the given option prices $C(t,x)$, namely
\begin{equation}\label{eq3.4}
 \frac{\sigma^2(t,x)}{2} = \frac{C_t(t,x)}{x^2 C_{xx}(t,x)} .
\end{equation}
Indeed, the Fokker-Planck equation implies -- at least on a formal level (cf.\ (\ref{Fokker0})) -- that $p_t(x)$ satisfies the PDE

\begin{equation}\label{Fokker}
\frac{\partial}{\partial t}p_t(x) = \frac{\partial^2}{\partial x^2}\left(\frac{\sigma^2(t,x)}{2}p_t(x)\right)
\end{equation}
Integrating in $x$, using (\ref{eq3.2}), and changing the order of derivatives quickly yields (\ref{eq3.4}).

We note that this beautiful argument is very much in line with Bachelier's reasoning in (\ref{lbhef}) and (\ref{heflbhef}) above pertaining to the case of constant volatility $\sigma$. We note in passing that Bachelier used instead of the wording ``volatility'' the more colorful term ``nervousness of the market''. 
%We also note that Bachelier's heuristic reasoning could formally be extended to the case of a diffusion coefficient $\sigma(t,x)$ depending on time and space yielding the more general Fokker-Planck equation (\ref{Fokker}) depending on $\sigma(\cdot,\cdot)$.

\vspace{1em}

%Indeed, a straightforward application of It\^o's formula to the martingale $C(t,S_t)$ yields

%\begin{equation}\label{eq3.5}
%dC(t,S_t) = C_t(t,S_t)dt +C_x(t,S_t) dS_t +\frac{1}{2}C_{xx}(t,S_t)\langle dS_t \rangle.
%\end{equation}
%Using (\ref{eq3.3}) to calculate the term $\langle dS_t \rangle$, and setting the drift term in (\ref{eq3.5}) to zero implies ``Dupire's formula'' (\ref{eq3.4}).

Of course, Dupire's formal arguments need proper regularity assumptions in order to be justified. There are two aspects: existence and uniqueness of the martingales fitting the given option prices $C(t,x)$. As regards the former, the question of existence amounts to a remarkable theorem by Kellerer \cite{Ke72, Ke73}: %\comment{MB: I changed this a bit since Kellerer makes no assertions on continuity.}
given a family $(\mu_t)_{0 \le t \le T}$ of probability distributions on $\R$ which is  increasing in convex order, there is a Markov martingale having these probabilities as marginals. %By weakly continuous we mean that, for $t_n \to t$, we have $\pi_{t_n}(f) \to \pi_t(f)$, for every bounded continuous function $f$ on $\R$. 
By increasing in convex order we mean that each $\mu_t$ has finite first moment and that $\mu_t (f)$ is nondecreasing in $t$, for every convex function $f$ on $\R$. 
Kellerer's theorem extends earlier work of Strassen \cite{St65} who establish the discrete time version of the result. We also note that the convex order condition on the marginal distributions is necessary as easily follows from Jensen's theorem.


% given a family $(\pi_t)_{0 \le t \le T}$ of probability distributions on $\R$ which is weakly continuous and increasing in convex order, there is a continuous martingale having these probabilities as marginals. By weakly continuous we mean that, for $t_n \to t$, we have $\pi_{t_n}(f) \to \pi_t(f)$, for every bounded continuous function $f$ on $\R$. By increasing in convex order we mean that each $\pi_t$ has finite first moment and that $\pi_t (f)$ is nondecreasing in $t$, for every convex function $f$ on $\R$.

Kellerer's theorem goes far beyond the simple formula (\ref{eq3.4}) and has been further refined, notably by Lowther in an impressive series of papers \cite{Lo08a, Lo08b, Lo08d}. We shall review these results in the subsequent sections.

However, from an application point of view, the existence question is not of primordial relevance. After all, the function $C(t,x)$ is an idealization of reality which has to be estimated from a finite set of given European option prices. In this context, it does not harm to make strong regularity assumptions on the smoothness and convexity (in the variable $x$) of the function $C(t,x)$  which justify the above argument. Under such assumptions Dupire's solution (\ref{eq3.4}) does make sense and the issue of existence is settled.

A different issue is the question of uniqueness. As we shall see below, this question is challenging and relevant -- at least from a mathematical point of view -- even in very regular settings, such as the Bachelier or the Black-Scholes model.

In order to formulate existence and uniqueness results for a process with given marginals, one has to specify the class of processes with respect to which we want to establish existence and uniqueness. Under proper regularity assumptions the unique solution should, of course, equal Dupire's solution. Dupire's process is a \textit{martingale with continuous paths, enjoying the Markov property}. Is Dupire's solution unique within this class? In a veritable tour de force Lowther  \cite{Lo08a, Lo08b} has shown  that the answer is yes, provided that we replace the word \textit{Markov} by the word \textit{strong Markov} and one restricts to continuous processes.

 We also refer to Theorem 6.1 in \cite{HiRoYo14} where a slightly different version of this theorem, credited to Pierre, is proved. These theorems settle the question of uniqueness in a very satisfactory way. We shall discuss Lowther's theorem in more detail in Section 6. 


To the best of our knowledge the following question remained open: is it really necessary to add the adjective \textit{strong} to the word \textit{Markov} in Lowther's uniqueness theorem? At least, if one is willing to accept strong regularity assumptions on the function $C(\cdot,\cdot)$ and the resulting process $S$ as defined in (\ref{eq3.3}) one may ask whether the Markov property alone is sufficient. We will focus on this question in the next section. % where we develop some ideas hinting in the direction that the answer is no. The construction of an explicit counter-example is quite technical and carried out in the twin paper \cite{BePaSc21a} where we show that, even for Gaussian marginals, there is a ``fake'' Brownian motion which is a continuous Markov martingale. Of course, this martingale must fail the strong Markov property in view of Lowther's uniqueness theorem.

 %In the accompanying paper \cite{BePaSc21a}  to the present one we shall show by means of an explicit example that even in the ``most regular'' case, i.e.\ even in Bachelier's model of Brownian motion, the answer is ``no''. Passing to a different terminology this means that there is a \textit{fake Brownian motion} which is a Markovian, continuous martingale that has Brownian marginals but differs from Brownian motion. 
 

 %We also show in [XXX] that this construction is of a generic feature and applies to a large class of given (marginals of) continuous Markovian martingales, e.g., the Black-Scholes model.

\section{An eye-opening example}

The subsequent example is known since the work by Dynkin and Jushkevich in the fifties  (see \cite{DyJu56}). 

\begin{example}\label{easy}

There is an $\R_+$-valued, continuous, Markov martingale which fails to be strongly Markovian.
\end{example}

\begin{proof}\label{easyproof}
We define the process $S=(S_t)_{0\le t\le 1}$ by starting at $S_0 = 1$ and subsequently proceeding in two steps. For $t \in [0,\frac{1}{2}]$ the process $S$ is a stopped geometric Brownian motion, i.e.,
\begin{equation} \label{geom}
S_t = \exp \left(B_{t \wedge \tau} -\frac{t \wedge \tau}{2} \right) \qquad \qquad 0\le t\le \frac{1}{2},
\end{equation}
where $B$ is a standard Brownian motion and $\tau$ is the first moment when $S$ hits the level $2$. 

For $\frac{1}{2} \le t \le1$, we distinguish two cases. If $S$ has been stopped, i.e., if $S_{\frac{1}{2}} = 2$, the process $S$ simply remains constant at the level $2$. If this is not the case, the process continues to follow a geometric Brownian motion, i.e.,
\begin{equation} \label{geom2}
S_t = \exp \left(B_t -\frac{t}{2} \right), \qquad \qquad \frac{1}{2} \le t \le 1 \quad \text{on the set} \quad  \left\{\tau > \frac{1}{2}\right\}  ,
\end{equation}

Obviously $S$ is a continuous martingale. The crucial feature is its Markovian nature: the Markov property follows from the fact that, for every fixed (deterministic) time $\frac{1}{2}\le t\le 1$, the probability for the geometric Brownian motion $S_t$  to be equal to $2$ is zero on the set $\{\tau>\frac{1}{2}\}$. Hence, for every fixed $\frac{1}{2} \le t \le 1$, the conditional law of $(S_u)_{t \le u \le 1}$ is \textit{almost surely} determined by the present value $S_t$ of the process.

Why does $S$ fail to be strongly Markovian? On the set $\{\tau > \frac{1}{2}\}$ define the stopping time $\vartheta$ as the first instance $u > \frac{1}{2}$ when $S_u$ equals the value $2$, which happens with positive probability during the interval $\left( \frac{1}{2} , 1 \right)$. The process $S$ therefore, at time $\vartheta \wedge 1$, takes the value $2$ on a non-negligible part of the set $\{\tau > \frac{1}{2}\}$. Of course, the random variable $S_{\tau}$ equals $2$ on the set $\{\tau\le \frac{1}{2}\}$ too. Hence there is no strong Markovian prescription for the process $S$ what to do after time $\vartheta$: without further information on the past the process $S$ cannot decide whether it should remain constant or continue to move on as a geometric Brownian motion.
\end{proof} 

Let us apply this example to the pricing of options of the form 
$(S_T -x)_+$. Fix $x \ge 2$. It is straightforward to calculate its price $P^{x}(t,z)$ at time $t$, conditionally on $S_t = z$, which is defined via
\begin{equation} \label{cond}
P^{x}(t,z) = \E [(S_T -x)_+|S_t=z],   \qquad \le t\le T, \, z \in \R.
\end{equation}
Letting $z=2$, we find
\begin{equation}\label{case2}
P^{x}(t,z) =0, \qquad 0 <t \le 1.
\end{equation}

On the other hand, for $z \neq 2$ and $\frac{1}{2} \le t <1$, the prices $P^{x}(t,z)$ are given by the usual Black-Scholes formula and therefore strictly positive. Hence, for $\frac{1}{2} \le t <1$, the option prices $z \to P^{x}(t,z)$ are  discontinuous at $z=2$. They also fail to be increasing and convex in the variable $x$ which a reasonable option pricing regime should certainly satisfy. On the other hand we note that these option prices -- strange as they might be -- do not violate the no arbitrage principle as they were legitimately derived from a martingale.

\vspace{1em}


The marginal probabilities of the process $S$ have an atom at the point $2$ which is rather unpleasant. One may ask whether it is possible to construct variants of the above example which have more regular marginals. 

Here is a straightforward modification: fix an uncountable compact $K$ in $\R_+$ with zero Lebesgue measure. For example one may take the classical Cantor set $K = \{1 + \sum_{n=1}^{\infty} \frac{\epsilon_n}{3^n} : \epsilon_n \in \{0,2\}\}$. We can modify the construction of Example \ref{easy} in three steps. For $0 \le t \le \frac{1}{3}$, let $(S_t)_{0 \le t \le \frac{1}{3}}$ be geometric Brownian motion starting at $S_0 =1$. For $\frac{1}{3} \le t \le \frac{2}{3}$, let $(S_t)_{ \frac{1}{3}\le t\le \frac{2}{3}}$ continue to be geometric Brownian motion, but stopped at the first hitting time $\tau$ of the compact set $K$. For $\frac{2}{3} \le t \le 1$ we again distinguish between the cases $\{\tau \le \frac{2}{3}\}$ and $\{\tau > \frac{2}{3}\}$. On the former set the process $S$ remains constant, i.e., $S_t = S_{\frac{2}{3}}$. On its complement $\{\tau > \frac{2}{3}\}$ the process $S$ continues to follow a geometric Brownian motion. The process $S$ thus enjoys all the features of Example \ref{easy} and, in addition, has continuous marginals. Note, however, that these marginals are not given by densities as they are not absolutely continuous with respect to Lebesgue measure.





\vspace{1em}

Turning back to the context of Example \ref{easy} there is another continuous Markovian martingale with the same marginals as $S$, inducing reasonable option prices. In fact, there is a continuous \textit{strong} Markovian martingale with this property and which is unique in this latter class (Theorem \ref{uniqueness theorem} below).

We only give an informal, verbal description of this strong Markov process. For $0 \le t \le \frac{1}{2}$ let $S$ be defined as in Example \ref{easy}. For $\frac{1}{2} \le t \le 1$ we again distinguish two cases. On the set $\{S_t \neq 2\} $ define $S$ to be geometric Brownian motion, but now we stop this motion when $S_t$ hits the value 2. On the other hand, on the set $\{S_t = 2\} $ the process $S$ starts an excursion from $S_t=2$ into the geometric Brownian motion (\ref{geom2}) with a certain intensity rate. We are free to choose this rate in such a way that the mass remaining at the atom $\{S_t=2\}$ equals precisely the constant mass which is prescribed by the given marginals of the process $S$.

We thus have indicared the construction of another continuous martingale having the same marginals as the process $S$ in Example \ref{easy}. One may check that the latter construction is strongly Markovian -- as opposed to the above construction in Example \ref{easy} -- and that the option prices are increasing and strictly convex in the variable $z$ as they should be. It will follow from Theorem \ref{uniqueness theorem} below that the latter martingale is the unique strong Markov solution for the given marginals.

%The argument in the preceding paragraph can also be applied to the case of the ``Cantor set construction'' above where the marginal distributions do not have atoms.


Note that this answers the question raised at the end of the last section: In Lowther's uniqueness theorem, it is not sufficient to consider Markovian (but not necessarily \emph{strongly} Markovian) martingales: as we have just seen, there exist two distinct continuous Markov martingales with the same 1-dimensional distributions. 

 In view of Dupiere's formula, this leads to the next question. It seems natural to conjecture that, provided the call prices are sufficiently regular in $t$ and $x$, there should be only one continuous Markov martingale matching these prices. Correspondingly one would ask:  can one obtain a similar example as above  with absolutely continuous (or even more regular) marginals?

 To the surprise of the present authors it turned out that the answer is ``yes'', even when we pass to the ``most regular'' situation when $S$ is a Brownian motion, i.e., in the Bachelier model (or the Black-Scholes model). The construction is more involved but resting on the above developed intuition, see the accompanying paper \cite{BePaSc21a}. %It will be done in section XXXX below where we also show that this construction applies to a large family of continuous Markovian martingales $S$, e.g., the Black-Scholes model.

%This simple example shows what can go wrong 


\section{Uniqueness of Dupire's diffusion}

There is a huge literature on one-dimensional processes inducing a given family of marginal distributions (see \cite{Ke72, MaYo02, HiRo12, HiRoYo14, BeHuSt16, Lo08b, HaKl07, FaHaKl15, Ho13, Ol08, Al08, BaDoYo11, KaTaTo15} among others). In particular, the late Marc Yor and his co-authors Hirsch, Profeta, and Roynette wrote the beautiful book \cite{HiPrRoYo11} on ``peacocks''.
  This is a  pun on the French acronym PCOC, for ``Processus Croissant pour l'Ordre Convexe'' and a \emph{peacock} is a stochastic process $(X_t)_{t\geq 0}$ for which the family of laws $ \text{law}(X_t), t\geq 0$ is increasing in the convex order. We take here the liberty to use the word peacock also for a family of probabilities $(\mu_t)_{t\geq 0}$ that increases in convex order.\footnote{Fortunately, ``Probabilit\'es Croissant pour l'Ordre Convexe'' still yields the same acronym.}

 To comply with this literature we find it more natural to pass from the multiplicative setting (\ref{eq3.3}) to the additive setting of a martingale diffusion

\begin{equation}\label{eq3.3add}
{dX_t}=\sigma(t,X_t)\, dW_t,\qquad 0\le t \le T.
\end{equation}
Hence we consider now processes taking possibly values in all of $\R$ and switch to the notation $X$ instead of the ``stock price'' $S$. We note, however, that this change is only for notational reasons, and everything below could also be done in the multiplicative setting of the previous sections.

Given a peacock $(\mu_t)_{t\geq 0}$ we may define option prices via
\begin{equation} \label{eq3.1ad}
C(t,x) = \E_{\mu_t}[(X_t-x)_+],
\end{equation}
where $\mu_t, t\geq 0$ denote the marginal probability measures and $x \in \R$. The `multiplicative' formula (\ref{eq3.4}) becomes in the additive setting
\begin{equation}\label{eq3.4add}
 \frac{\sigma^2(t,x)}{2} =  \frac{C_t(t,x)}{ C_{xx}(t,x)} .
\end{equation}

We can now cite Lowther's complete solution of the uniqueness problem within the class of continuous, strong Markov martingales. We stress (and admire) that this theorem does not require any additional regularity assumptions \cite[Theorem 1.2]{Lo08a}.

\begin{theorem}[Lowther] \label{uniqueness theorem}
Let $X= (X_t)_{0\le t\le1}$ and $Y= (Y_t)_{0\le t\le1}$ be  $\R$-valued, {\it continuous, strong Markov martingales}. 
If  $X$ and $Y$ have the same marginal distributions they also have the same joint distributions. 
\end{theorem}

The proof of this theorem is highly technical and its presentation goes far beyond the scope of the present paper. Instead, we formulate a ``toy'' version of the theorem under strong regularity assumptions. We then analyze why the notion of \textit{strong} Markovianity is key in the above theorem and finally give some hints on the strategy of the proof of Theorem \ref{uniqueness theorem}.

\vspace{1em}

Here is the ``toy'' version of this theorem imposing strong regularity assumptions which make life easier:

\begin{assumption}\label{Ass1}
 We suppose that the process $X$ is given by $X_0 = x_0$ and (\ref{eq3.3add}), where $\sigma(t,x)$ is sufficiently smooth to guarantee that there is a unique strong solution $X$. We also suppose that $X_T$ has finite second moment and that the function $C(t,x)$ is strictly convex in the variable $x$, strictly increasing in the variable $t$, and satisfies standard It\^o smoothness assumptions, i.e., twice continuously differentiable in $x$ and once continuously differentiable in $t$. We also assume that, for every $x \in \R$, the pricing function $(t,z) \to P^{X,x}(t,z)$ defined via

\begin{equation} \label{condadd}
P^{X,x}(t,z) = \E [(X_T -x)_+|X_t=z],   \qquad 0 \le t\le T, \, z \in \R.
\end{equation} 
also satisfies these standard It\^o assumptions.

\end{assumption}

These assumptions are strong enough to guarantee that the function $C(t,x)$ indeed satisfies Equation (\ref{eq3.4add}). 

%\begin{equation}\label{Kolm}

%\end{equation}

\begin{theorem}\label{baby}
Let $X=(X_t)_{0\le t\le T}$ satisfy Assumption \ref{Ass1}.
Let $Y=(Y_t)_{0\le t\le T}$ be another continuous Markov martingale such that $X_t$ and $Y_t$ have the same distribution, for every $0 \le t \le  T$.
For fixed strike price $x \in \R$, let  $P^{Y,x}(t,z)$ be the corresponding option prices defined via
\begin{equation} \label{condadd'}
P^{Y,x}(t,z) = \E [(Y_T -x)_+|Y_t=z],   \qquad 0 \le t\le T, \, z \in \R,
\end{equation}
and {\bf assume} that, for every $x$, the functions $(t,z) \to P^{Y,x}(t,z)$ also satisfy the above standard It\^o smoothness assumptions.
Then $P^{X,x}(t,z) = P^{Y,x}(t,z)$, for all $t,x,z$ and the processes $X$ and $Y$ have the same joint distributions.
\end{theorem}

\begin{proof}
As the function $(t,z) \to P^{Y,x}(t,z)$ is {\bf assumed} to satisfy the standard It\^o conditions we may apply It\^o's  formula to obtain
\begin{equation}\label{Ito}
dP^{Y,x}(t,Y_t) = P_t^{Y,x}(t,Y_t) \, dt  + P_z^{Y,x}(t,Y_t)\, dY_t + \frac{1}{2} P_{zz}^{Y,x}(t,Y_t)\, d\langle Y {\rangle}_t,
\end{equation}
where $\langle Y {\rangle}_t$ denotes the quadratic variation process of the continuous, square integrable martingale $Y$ (see 
\cite[Theorem 7.6.4]{Ku06}).
By \eqref{condadd'} the process $(P^{Y,x}(t,Y_t))_{0 \le t \le 1}$ is a martingale. The martingale condition implies that the drift term vanishes so that the equality $\frac{1}{2} P_{zz}^{Y,x}(t,Y_t) \, d\langle Y {\rangle}_t = - P_t^{Y,x}(t,Y_t)\, dt$ holds true in the following sense:
\begin{equation}\label{mart}
\frac{1}{2} \int_A  P_{zz}^{Y,x}(t,Y_t)\, d\langle Y {\rangle}_t = - \int_A P_t^{Y,x}(t,Y_t)\, dt,
\end{equation}
whenever we integrate over a predictable set $A \subseteq [0,T] \times \mathcal C[0,T]$. In particular, the finite measure $d\langle Y {\rangle}_t$ on the predictable sigma-algebra of subsets of $[0,T] \times \mathcal C[0,T]$ has to be absolutely continuous with respect to the measure $d\langle X {\rangle}_t = \sigma^2(t,\omega_t)\, dt$. Here we have used the assumption that $t \to C(t,z)$ is strictly increasing in $t$ so  that $\frac{\sigma^2(t,z)}{2} =  \frac{C_t(t,z)}{ C_{zz}(t,z)}$ is strictly positive. We thus may define the function 
\[ \frac{\rho^2 (t,z)}{2} := \frac{P_t^{Y,x}(t,z)}{P_{zz}^{Y,x}(t,z)} \]
so that 
$d\langle Y {\rangle}_t =  \rho^2 (t,Y_t) \, dt$. 
We therefore must have that $Y$ may be represented as in (\ref{eq3.3add}), with $\sigma$ replaced by $\rho$. This implies the relation
\begin{equation}\label{forw}
 C_t(t,x) = \frac{\rho^2(t,x)}{2} { C_{xx}(t,x)} .
 \end{equation}
Comparing with (\ref{eq3.4add}) we obtain $\rho^2 = \sigma^2$ which shows the identity of $X$ and $Y$ in distribution.
\end{proof}




The above theorem provides a sufficient set of regularity assumptions to substantiate the statement in Dupire's paper \cite{Du94}: ``...we can recover, up to technical regularity  assumptions,  a  unique  diffusion  process''.

Of course, one could do some massaging of the above argument to somewhat weaken the very strong Assumptions \ref{Ass1} which we have imposed. But there is a long and thorny road, going far beyond simple cosmetic changes, to arrive at Lowther's  Theorem \ref{uniqueness theorem}.

In Theorem \ref{baby} the strong regularity assumptions implied in particular the strong Markov property of the process $X$. We stress once more that, in the setting of Lowther's theorem \ref{uniqueness theorem}, the \textit{strong} Markov property is the key assumption.
\vspace{1em}

Passing to Lowther's notation and looking at \eqref{condadd}, a crucial step in the above argument is to start from a convex, increasing, and Lipschitz-one function $g(z)$, such as $g(z)=(z-x)_+$, to its conditional expectations 
\begin{equation} \label{conda}
f(t,z) = \E [g(Y_T)|Y_t=z],   \qquad 0 \le t\le T,\, z \in \R.
\end{equation}
In order to start a chain of arguments one has to verify that $f(t,z)$ is a ``nice'' function.
When looking at Example \ref{easy} and its variants we have seen that in this case, for $g(z)=(z-x)_+$, this is not at all the case. 
Its conditional expectation $f(t,z)$ lacked each of the following desired properties: continuity, monotonicity and convexity in $x$. 

Contrary to this lamentable breakdown of regularity, we shall verify in Corollary \ref{corr} that the strong Markov property guarantees that the following three properties  are inherited from $g(\cdot)$ to each $f(t,\cdot)$: convexity, monotonicity, and $1$-Lipschitz continuity (which serves as a more quantitative version of continuity). This preservation of regularity is a decisive feature of Lowther's proof.


%
%Let us turn to Lowther's strategy in [].
%
%He started fromthe symmetric relation between the ``backward solution'' $g(\cdot,\cdot)$ and the ``forward solution'' $C(\cdot,\cdot)$ 
%
%Ist die Bezeichnungsweise richtig?
%
%\begin{equation}\label{Lowther1.6}
%\frac {\partial f}{\partial t}  \frac {\partial^2 C}{\partial x^2} + \frac {\partial C}{\partial t}\frac {\partial^2 f}{\partial x^2} =0.
%\end{equation}
%which we have applied in the above proof to $f(t,z) = P^{Y,x}(t,z)$.
%
%In our argument we have made sure that the above functions are well behaved so that (\ref{Lowther1.6}) makes good sense. But in general .....
%



\section{Coupling strong Markov processes}

What is the salient property which distinguishes the strong Markov property from the Markov property in our context? While the former condition allows for Lowther's uniqueness theorem to holde true, in Example \ref{easy} we have seen that there may be different continuous Markov martingales inducing the same marginals. The following well-known concept is the key to understanding the difference.



\begin{definition}\label{monotone}
For probability measures $\pi^1$ and $\pi^2$ on $ \R$, we say that $\pi^2$ dominates $\pi^1$ of first order if, for every $a \in \R$, we have $\ \pi^2[a, \infty[\, \ge \pi^1[a, \infty[$.
\end{definition}
 
 We shall show in the next proposition that the strong Markov property of a continuous martingale implies that the transition probabilities $(\pi_x^{s,t})_{x\in \R}$
 \begin{equation}\label{condpro}
 \pi_x^{s,t}[A] = \P [X_t \in A | X_s = x],
 \end{equation}
where $s<t$ and $A$ is a Borel set in $\R$, are increasing of first order in the variable $x$, for every $s<t$.
%
%\comment{MB: Ich wuerde den Text in cyan gerne weglassen.}
%{\ccyan
%\vspace{1em}
%
%Hirsch, Roynette, and Yor \cite{HiRoYo14} provided the following nice argument for this on an infinitesimal level, very much in the spirit of Bachelier's arguments. It makes precise -- under strong regularity -- conditions the idea that ``paths do not cross''.
%
%\vspace{1em}
%
%XXXX   Ich habe die Passage ins Tex file kopiert, die man noch etwas umformulieren muss. Ich wuerde auch hier (noch) nicht von Lipschitz-Markov reden. XXXXX
%  \vspace{1em}
%
%
%The proposition allows us to visualize very nicely the paths $X^x_t$ and $X^{x+dx}$, starting from two infinitesimal neighbours $x$ and $x+dx$, in the setting of two stochastic integrals driven by the same Brownian motion. The infinitesimal distance of their positions may very well change during their journeys, but it cannot become zero. In addition, if we make sure that the above encountered local martingale is a true martingale (which can easily be done by adding one more regularity assumption), we see that \textit{in average} the infinitesimal distance remains constant.
%
%\vspace{1em}
%
%After this motivating result} 

We follow Hobson who applied a well-known technique, namely the ``joys of coupling'' (to quote his paper \cite{Ho98c}) in the present context.

\begin{proposition}\label{Hobson}
Let $X = (X_t)_{0 \le t \le T}$ be a continuous strong Markov process with transition probabilities $ \pi_x^{s,t}[ \cdot]$. Then, for $0 \le s \le t \le T$, and for $x < y$ the probability $\pi_y^{s,t}$ dominates  $\pi_x^{s,t}$ of first order.
\end{proposition}

\begin{proof}
Fix $s, t$ and $ x<y $ as above and let $(X^x_u)_{s \le u \le t}$ and $(X^y_u)_{s \le u \le t}$ be independent copies of the process $X$, starting at $X^x_s=x$ and $X^y_s = y$, both defined on the same filtered probability base. Define the stopping time $\tau$ as the first moment $u$ when the process $X^x_u$ equals the process $X^y_u$, if this happens for some $u \in [s,t]$; otherwise we let $\tau = \infty$.

Define the process $\tilde{X}^x$ by
\[
	\tilde{X}^x_u = 
	\begin{cases} 
		X_u^x & \text{for }s \leq u \leq \tau, \\
		X_u^y & \text{for }\tau < u \leq t.
	\end{cases}	
\]
We clearly have 
\begin{equation}\label{ineq}
\tilde{X^x_t} \le X^y_t,\qquad \text{for all } 0 \le t \le T.
\end{equation}
Indeed, if $\tau = \infty$, the paths of $(X^x_u)_{s \le u \le t} = (\tilde{X}_u^x)_{s \le u \le t}$ and $(X^y_u)_{s \le u \le t}$ never touch, so that we even have a strict inequality by continuity of the processes. If $\tau < \infty$, then $\tilde{X}^x$ and ${X^y}$ have ``joined'' at time $\tau$, and follow the same trajectory. Hence $\tilde{X}_u^x = X^y_u$, for $\tau \le u \le t$.

Inequality \ref{ineq} implies that the law of $X^y_t$ dominates the law of $\tilde{X}_t^x$ in first order.
\end{proof}

\begin{corollary}\label{corr}
Let $X = (X_t)_{0 \le t \le T}$ be a continuous strong Markov process with marginal laws $\mu_t$ and transition probabilities $\pi_x^{s,t}[ \cdot]$. Let $0 \le s \le t \le T$  and $z \mapsto g(z)$ be a measurable $\mu_t$-integrable function, and define the conditional expectation similarly as in (\ref{conda})

\begin{equation} \label {condad}
f(s,x) = \E [g(X_t)|X_s=x],   \qquad  x \in \R.
\end{equation}
Then the following assertions hold true.

(i) If $z \mapsto g(z)$ is increasing, then so is $x \mapsto f(s,x)$, for every $0 \le s \le t \le T$.

\noindent If, in addition, we assume that $X$ is a martingale, we also have the following two assertions.

(ii) If $z \mapsto g(z)$ is 1-Lipschitz , then so is $x \mapsto f(s,x)$, for every $0 \le s \le t \le T$.

(iii) If $z \mapsto g(z)$ is convex, then so is $x \mapsto f(s,x)$, for every $0 \le s \le t \le T$.



\end{corollary}

\begin{proof}
Assertion (i): this is just a reformulation of Proposition \ref{Hobson}.

Assertion (ii): If $g$ is 1-Lipschitz, then $g(x) +x$ is increasing. As $X$ is a martingale we have $\E [X_t|X_s=x] = x$. By (i) $f(s,x) + x$ is increasing. By the same token $f(s,x) - x$ is decreasing which readily shows that $x \to f(s,x)$ is Lipschitz.

Assertion (iii): We follow the proof of \cite[Theorem 3.1]{Ho98c}.
For convex $g$ and $x < y < z$ we have to show that
\begin{equation}\label{convex}
(z-x) f(s,y) \le (z-y)f(s,x) + (y-x) f(s,z).
\end{equation}

Fix $x < y < z$ and choose three independent copies $X^x, X^y,X^z$ of the process $X$, starting at time $s$ from the initial values $x,y$, and $z$. To simplify notation we denote the resulting triple of processes $\left(  {X^x},X^y,{X^z} \right)$ by $(X,Y,Z)$. We define coupling times
similarly as above. Let $\tau^x$ be the first moment $u > s$ when $X(u) = Y(u)$; similarly $\tau^z$ is defined as the first moment when $Y$ and $Z$ meet. Finally, let $ \tau = \tau^x \wedge \tau^z \wedge t$. 
%As in the proof of Proposition \ref{Hobson} we change the process $X^x$ to follow $X^y$ after $\tau^x$ and $X^z$ to follow $X^y$ after $\tau^z$. 
This time we leave the processes unchanged; rather we shall argue on the three disjoint (up to null sets) sets $\{\tau = \tau^x\}$, $\{\tau = \tau^z\}$, and $\{\tau = t\}$.

We start with the latter on which we have $X_t < Y_t < Z_t$.
By the convexity of $g$ we have
\begin{equation}\label{convexrandom}
(Z_t-X_t) g(Y_t) \le (Z_t-Y_t) g(X_t) + (Y_t-X_t) g(Z_t).
\end{equation}
so that
\begin{equation}\label{convexexp}
\E\left[(Z_t-X_t) g(Y_t) - \left( (Z_t-Y_t) g(X_t) + (Y_t-X_t)g(Z_t)\right);\tau = t\right] \le 0.
\end{equation}

On $\{\tau = \tau^x\}$ we have $X_t = Y_t $ so that the last term in inequality (\ref{convexrandom}) vanishes. Also the first and the middle term are equal so that (\ref{convexrandom}) holds true (with equality) on the set $\{\tau = \tau^x\}$. In particular,
\begin{equation}\label{convexexp1}
\E\left[(Z_t-X_t) g(Y_t) - \left( (Z_t-Y_t) g(X_t) + (Y_t-X_t)g(Z_t)\right);\tau = \tau^x\right] \le 0.
\end{equation}
The same reasoning applies to $\{\tau = \tau^z\}$. 
\begin{equation}\label{convexexp2}
\E\left[(Z_t-X_t) g(Y_t) - \left( (Z_t-Y_t) g(X_t) + (Y_t-X_t)g(Z_t)\right);\tau = \tau^z\right] \le 0.
\end{equation}
Summing over these three sets we obtain
\begin{equation}\label{convexexp3}
\E\left[(Z_t-X_t) g(Y_t) - \left( (Z_t-Y_t) g(X_t) + (Y_t-X_t)g(Z_t)\right)\right] \le 0.
\end{equation}
Finally we use independence and the martingality of $X,Y$, and $Z$ to obtain
\begin{equation}\label{convex1}
(z-x) \E[g(Y_t)|Y_s=y] \le  (z-y)\E[g(X_T|X_s=x] + (y-x)\E[g(Z_t)|Z_s=z],
\end{equation}
which is tantamount to (\ref{convex}).
\end{proof}

%\begin{remark}\label{remark1} Property $(ii)$ goes back to Kellerer's work \cite{Ke72}. It has been called ``Lipschitz Markov'' property in \cite{HiRoYo14} and simply ``Lipschitz'' property in \cite{Lo09}. As we have seen it quickly follows from the first order dominance (Proposition \ref{Hobson}) and the martingale property of $X$.
%\end{remark} 

We can reformulate the message of Corollary \ref{corr} (ii) in the spirit of Bachelier by considering the Wasserstein cost $\mathcal{W}_1 (\pi_x^{s,t}[ \cdot] , \pi_y^{s,t}[ \cdot])$ of the horizontal transport of the conditional probability measures $ \pi_x^{s,t}[ \cdot]$ to $ \pi_y^{s,t}[ \cdot]$.

Recall that, for probabilities $\mu, \nu$ on the real line the Wasserstein-1 distance is given by 
\[ \mathcal{W}_1(\mu, \nu):=\inf_{\pi \in \text{cpl}(\mu, \nu)} \int |x-y|\, d\pi(x,y), \]
where $\text{cpl}(\mu, \nu)$ denotes the set of all probabilities on $\R^2$ having $\mu, \nu$ as marginal measures, see e.g.\ \cite{Vi09} for an extensive overview of the field of optimal transport. 

\begin{definition}
Let $\pi$ be a probability on  $\R^2 $ and write $\mu$ for its projection onto the first coordinate and  $(\pi_x)_x$ for the respective disintegration so that $\pi = \int_\R \pi_x d\mu(x)$. Then $\pi$ is called a Lipschitz-kernel if for all $x, y $ in a set $X$ with  $\mu(X)=1$ 
\begin{align}\label{kLip}
\mathcal W_1(\pi_x, \pi_{y})\leq |x-y|.
\end{align} 
\end{definition} 
We call $\pi$ a \emph{martingale coupling} if $\int y\, d\pi_x(y)=x$, $\mu$-a.s. It is then straight-forward to see that for a martingale coupling $\pi$ the following are equivalent: 
\begin{enumerate}
\item $\pi$ is a Lipschitz-kernel.
\item for all $x, y $ in a set $X$ with $\mu(X)=1$ we have
$
\mathcal W_1(\pi_x, \pi_{y})= |x-y|.
$
\item for all $x, y$ in a set $X$ with $\mu(X)=1$, $x\leq y$ the measure $\pi_x$ is dominated in first order by $\pi_{y}$.  
\end{enumerate}

\begin{definition}\label{Wass} 
Let $X$ be an $\R$-valued Markovian process. Then X has the Markov Lipschitz property if for all $s \leq t$, the law of $(X_s, X_t)$ is a Lipschitz kernel. 
\end{definition}
To give yet another characterisation of Lipschitz-Markov processes, recall that a process $X$ is Markov if and only if for all $s\leq t$ and every bounded measurable function $f$ there is a measurable function $g$ such that 
$$\E[f(X_s)|\mathcal F_s]= g(X_t).$$ 
A process $X$ is Lipschitz-Markov if and only if for all $s\leq t$ and every 1-Lipschitz function $f$ there is a 1-Lipschitz function $g$ such that 
$$\E[f(X_s)|\mathcal F_s]= g(X_t).$$
This is a straightforward consequence of the Kantorovich-Rubinstein theorem that provides a dual characterization of the Wasserstein-1 distance through 1-Lipschitz functions. 

We now can resume the crucial role of the strong Markov property

\begin{corollary}\label{SLM}
Let $M$ be a continuous Markov martingale. Then $M$ is Lipschitz-Markov if and only if it is strong Markov. 
\end{corollary}
\begin{proof} 
By Proposition \ref{Hobson} / Corollary \ref{corr} every strong Markov-martingale is Lipschitz Markov.  
That a Lipschitz Markov martingale is strongly Markov is proved in the same way as one establishes the strong Markov property for Feller processes. See e.g.\ \cite[Theorem 1.68]{Li10}. \end{proof}

To the best of our knowledge, Lipschitz kernels play a crucial role in all known proofs of Kellerer's theorem. The decisive property is the following: 
\begin{proposition}\label{LMclosed}
Consider the space of $\mathcal P(\mathcal D[0,1])$ of probability measures on the Skorokhod space equipped with the convergence of finite dimensional distributions.  Then the set of Lipschitz-Markov martingales is closed. 
\end{proposition}
In contrast, the set of Markov-martingales is not closed. 
See e.g.\ \cite{BeHuSt16} for the (simple) proof of Proposition \ref{LMclosed}.


\section{Continuity of the martingale solution}

An important question in the present context is the following: under which conditions on a peacock $(\mu_t)_{0\le t \le 1}$, as defined in Section 4 above, is there a 
strong Markov martingale with \textit{continuous trajectories} having the given marginals. We only focus on the one-dimensional case and mention that the corresponding question for higher dimensions remains wide open. The one-dimensional case, however, is fully understood by now, again by the definitive work of Lowther \cite[Theorem 1.3]{Lo08b}:
\begin{theorem}[Lowther]%\label{LowThm}
Let $(\mu_t)_{t\geq 0}$ be a peacock and assume that $t\mapsto \mu_t$ is weakly continuous and that each $\mu_t$ has convex support. Then there exists a unique strongly continuous Markov martingale $X$ such that such that $X_t\sim \mu_t, t\geq 0$. 
\end{theorem}


%\comment{should we also cite other people here}

%\comment{WS Wir m\"ussen besprechen, ob wir auf SBM abstellen oder sagen, dass wir nur irgendeine stetige Lipschits Markov Lösung für mu und nu brauchen}



We do not show Lowther's theorem in full generality, but again we want to isolate a sufficient set of assumptions that allows us to present a (comparably simple) self-contained proof of the continuity theorem. 

\begin{remark}\label{LMkernelE}
A key ingredient of the proof is that for  probabilities $\mu, \nu$ in convex order there exists a continuous martingale $(X_t)_{0 \le t \le 1}$, $X_0 \sim \mu$, $X_1\sim \nu$  which is strongly Markovian (and hence Lipschitz Markov). For instance we can take $X$ to be a \emph{stretched Brownian motion}, that is, a solution to a continuous time martingale transport problem, see \cite{BaBeHuKa20}. Another possibility would be to apply an appropriate deterministic time change to Root's solution \cite{Ro69} (see \cite{CoWa12} for the case of a non-trivial starting distribution) of the Skorokhod embedding problem. We note that the martingale transport approach is also applicable to measures $\mu, \nu$ defined on $\R^d, d>1$. This could be interesting in view of a possible multi-dimensional  extension of Lowther's theorem but this is not within the scope of the present article.
\end{remark}




\begin{assumption}\label{Ass}

Let $ (\mu_t)_{0 \le t \le 1}$ be a one-dimensional peacock centered at zero, with densities $p_t(x)$ and finite second moments $m^2_2(\mu_t) = \int_{-\infty}^\infty x^2 \, d\mu_t(x) = \int_{-\infty}^\infty x^2 p_t(x)\, dx$ such that the function $ t \to m^2_2(\mu_t)$ is continuous. 

We assume that there is a -- bounded or unbounded -- open interval $I \subset \R$ supporting each $\mu_t$ such that, for each compact subset $ K \subset I$, the Lebesgue densities $p_t(x)$ of $\mu_t$ are bounded away from zero, uniformly in $x \in K$ and $t \in [0,1]$.

%We also assume that, for every strike price $k \in I$, the price of the option  $ x \to  \E_{\mu_t}[(x - k)_+]=\int_k^{\infty}xd\mu_t(x)$ is strictly increasing in $t \in [0,1]$.


%$0 \le t_0 < t_1 \le 1$, the pair $(\mu_{t_0},\mu_{t_1})$ is irreducible in the sense of \cite{BaBeHuKa20}.
\end{assumption}



It will be convenient to suppose (w.l.g.\ via a deterministic time change) that $ t \to m^2_2(\mu_t)$ is affine. More precisely we may assume that $m^2_2(\mu_{t+h}) - m^2_2(\mu_t)= h$ so that for all martingales $M$ with $law( M_t) = \mu_t$ and $law( M_{t+h} )= \mu_{t+h}$

\begin{equation} \label{1/n}
\E \left[\left(M_{t+h} - M_t \right)^2\right] = h.
\end{equation}


\begin{theorem}\label{cont}
Under the above assumptions there is a \underline{continuous} strong Markov martingale $M = (M_t)_{0 \le t \le 1}$ with given marginals $ (\mu_t)_{0 \le t \le 1}$.
\end{theorem}

There is an obvious and well-known strategy for the proof. We want to obtain the desired martingale $M$ as a limit of approximations which fit the peacock $(\mu_t)_{0 \le t \le 1}$ on finitely many points of time. As in \cite{HiRoYo14} it is convenient to do so along the ordered set $\cal S$ of finite partitions $ S = \{s_0,\dots,s_n\} $ of $[0,1]$, where $0 = s_0 < \dots < s_n =1$.



 For each $S \in \cal S$ we choose a continuous strong Markov martingale $M^S$ having the given marginals at each time $s_i \in S$, the existence of  $M^S$  is a direct consequence of Remark \ref{LMkernelE}. %standard by now and there are suitable constructions in the literature: E.g.\ given $s_i, s_{i+1}$, one can define $M^S$ to be the stretched Brownian motion (defined in \cite{BaBeHuKa20}) between the marginals $\mu_{s_i}, \mu_{s_{i+1}}$ and then concatenate the resulting processes. Another solution would be to apply an appropriate deterministic time change to Root's solution \cite{Ro69} of the Skorokhod embedding problem.
 
Identifying the martingales $M^S$ with their induced measures on the path space $\mathcal C[0,1]$ this family of measures is tight, if considered on Skorohod space $\mathcal D[0,1]$, equipped with the topology of convergence of finite dimensional distributions.
Hence we can find a cluster point $M$ in the set $\mathcal P(\mathcal D[0,1])$ of probability measures on $\mathcal D[0,1]$, see e.g.\ \cite{BeHuSt16} for the straightforward argument.
By refining the filter $\mathcal S$ we may suppose that $M$ is a limit point. 

We fix such a limiting process $M$ which by  Proposition \ref{LMclosed} is a Lipschitz Markov martingale.
These arguments again are standard by now and, e.g., well presented in the papers \cite{HiRoYo14, BeHuSt16}.
A priori the martingale $M$ has c\`adl\`ag trajectories.
Our present task is to show the \textit{continuity} of the trajectories of the limiting process $M$ under the above assumptions.

We first give a general criterion for the continuity of a limiting martingale $M$ which is somewhat reminiscent of the classical Kolmogorov continuity criterion \cite[Theorem 1.8]{ReYo99}.

\begin{proposition}\label{criterion}
	Let $(M^i)_{i \in I}$ be a net of $\R$-valued continuous strong Markov martingales  $M^i = (M^i)_{0 \le t \le 1}$ and $M$ its limit in the set of probabilities on Skorohod space $\mathcal P(\mathcal D[0,1])$ with respect to convergence in finite dimensional distribution.
	Suppose that there are constants $C_1>0$ and $\beta > 0$ such that, for every $0 \le t_0 < t_0 + h \le 1$ and every $i \in I$
	\begin {equation}\label{BMO}
	\lVert (M^i_t)_{t_0 \le t \le t_0 + h}\rVert_{{BMO}_1} := \sup_\tau\{ \lVert \E[|M^i_{t_0 + h} -M^i_\tau| \mid {\cal{F}}_{\tau}] \rVert_\infty   \}   \le  C_1 h^\beta.
	\end{equation}
	Then the martingale $M$ has continuous trajectories.
\end{proposition}

In (\ref{BMO}) $\tau$ runs through the $[t_0,t_0 + h]$-valued stopping times with respect to the natural filtration of $ (M^i_t)_{t_0 \le t \le t_0 + h}$. 
As $M^i$ is strong Markov, condition (\ref{BMO}) is tantamount to the requirement that the first moments $m_1(\pi_x^{i,\tau ,t_0 +h}) = \int_{-\infty}^\infty |y-x| \, d\pi_x^{i ,\tau ,t_0 +h}(y)$ of the transition probabilities $\pi_x^{i,\tau,t_0+h}$ from $M^i_\tau = x$ to $M^i_{t_0 + h}$ satisfy 
\begin{equation}\label {beta}
 	m_1( \pi_x^{i,\tau,t_0 +h})  = \int_{-\infty}^\infty |y-x| \, d\pi_x^{i ,\tau ,t_0 +h}(y) \le C_1 h^\beta, 
\end{equation}
for $\mu^i_\tau$-almost all $x \in \R$, where $\mu^i_\tau$ denotes the law of $M^i _\tau$.

An important feature of the $BMO$ norms for \textit{continuous} martingales is that, by  the John-Nirenberg inequality, all $BMO_q$ norms are equivalent, for $1 \le q < \infty$ (see, e.g.\ \cite[Corollary 2.1]{Ka06b}). Applying this fact to the present context, inequality (\ref{BMO}) is equivalent to the existence of a constant $C_q > 0$ (for some or, equivalently, for all $1 \le q < \infty$) such that

\begin {equation}\label{BMOq}
\lVert (M^i_t)_{t_0 \le t \le t_0 + h}\rVert_{\text{BMO}_q} := \sup_\tau\left\{ \left\lVert \E[|M^i_{t_0 + h} -M^i_\tau|^q \mid {\cal{F}}_{\tau}]^{\frac{1}{q}} \right\rVert_\infty \right\} \le C_q h^\beta.
\end{equation}
% Again this can be rephrased as 
% \begin{equation}\label{betaq}
%  	m_q(\pi_x^{i ,t_0 ,t_0 +h}) = \left(\int_{-\infty}^\infty |y-x|^q \, d \pi_x^{i ,t_0 ,t_0 +h}(y)\right)^{\frac{1}{q}} \le C_q h^\beta. 
% \end{equation}

\begin{proof}[Proof of Proposition \ref{criterion}]
Suppose that $M$ fails to be continuous and let us work towards a contradiction to \eqref{BMOq} when $q > \frac{1}{\beta}$.
Assume $M$ has jumps of size bigger than $3a > 0$ with probability bigger than $\kappa >0$, i.e.
\[
	\mathbb P\left( \left\{ \exists t \in [0,1] \colon |M_{t} - M_{t-}| \geq 3a \right\} \right) > \kappa.
\]
As $M$ has c\`adl\`ag paths, there is $h_0 > 0$ such that for all $0 < h \leq h_0$
\[
	\mathbb P\left( \left\{ \exists k \in \mathbb N \colon |M_{(kh)\wedge 1} - M_{((k-1)h)\wedge 1}| \ge 2a \right\} \right) > \kappa.
\]
By the pigeon hole principle, we can find for each $0 < h \leq h_0$ a time $t_0 \in [0,1]$ with
\begin{equation}\label{a}
	\P\left( \left\{ |M_{(t_0 +h)\wedge 1} - M_{t_0}| \ge 2a \right\}  \right) > h\kappa.
\end{equation}
In view of the convergence in finite dimensional distributions of $(M^i)_{i \in I}$ to $M$, we find for each $0 < h \leq h_0$ a time $t_0 \in [0,1]$ (w.l.g.\ $t_0 + h \leq 1$) and an index $i \in I$ with
\begin{equation}\label{a'}
	\P\left( \left\{ |M^i_{t_0 + h} - M^i_{t_0}| \ge a \right\} \right) \geq h \kappa.
\end{equation}
Fixing such an index $i \in I$, it follows that there is a set $A \subset \R$ of positive measure with respect to the law of $M^i_{t_0}$ such that, for all $x \in A$,
\begin{equation}\label{acond}
	\P\left( \left\{ |M^i_{t_0 +h} - M^i_{t_0}| \ge a \mid M^i_{t_0} = x \right\} \right) \geq h\kappa.
\end{equation}
For $x \in A$ we therefore have
\begin{equation}\label{mq}
	m_q(\pi_x^{i ,t_0 ,t_0 +h}) = \left(\int_{-\infty}^\infty |y-x|^q \, d\pi_x^{i,t_0 ,t_0 +h}(y)\right)^{\frac{1}{q}} \geq a(h\kappa )^{\frac{1}{q}}.
\end{equation}
As $q > \frac{1}{\beta}$ we can choose $0 < h \leq h_0$ sufficiently small, with $a (\kappa h)^{\frac{1}{q}} > C_q h^\beta$, and arrive at the desired contradiction to \eqref{BMOq}
\begin{equation*}
	C_q h^\beta \geq \lVert (M_t^i)_{t_0 \leq t \leq t_0 + h} \rVert_{BMO_q} \geq m_q(\pi_x^{i,t_0,t_0+h}) \geq a (h\kappa)^{\frac{1}{q}} > C_qh^\beta.
\end{equation*}
\end{proof}


Turning back to the above defined family $(M^S)_{S \in \cal S}$ of martingales we shall establish an inequality of the type (\ref{beta}) for the transition probabilities $\pi_x^{S ,t_0 ,t_0 +h}$, using the fact that the functions $x \mapsto \pi_x^{S,t_0,t_0+h} $ are Lipschitz-Markov.

\begin{lemma}\label{sublemma}
	Let $(\mu_t)_{0 \le t \le 1}$ be a peacock satisfying Assumption \ref{Ass}.
	Fix a compact set $K \subset I$.
	For all $h > 0$ sufficiently small and $x \in K$, there is a constant $D > 0$ such that for all $S \in \mathcal S$ and $0 \leq t_0 \leq t_0 + h \leq 1$, with $t_0, t_0 + h \in \mathcal S$, the first moments of the transition measures $\pi_x^{S ,t_0 ,t_0 +h}$ can be estimated by
	\begin{equation} \label{m1}
		m_1\biggl(\pi_x^{S ,t_0 ,t_0 +h}\biggr) := \E_\P\left[\left|M^S_{t_0 +h} - M^S_{t_0}\right|  \biggm| M^S_{t_0}=x\right] 
		=\int_{-\infty}^\infty \lvert y-x\rvert \, d\pi_x^{S ,t_0 ,t_0 +h}(y)
		\le D h^{\frac{1}{4}},
	\end{equation}
	where $\P$ denotes the law of the martingale $M^S$.
\end{lemma}

\begin{proof}
We first suppose that $I=\R$. By \eqref{1/n} and Jensen's inequality we have
\[
	\E_\P[|M^S_{t_0+h} - M^S_{t_0}|] \le  h^{\frac{1}{2}}.
\]
 We may rewrite this inequality in the form 
\begin{align} \nonumber
\E_\P[|M^S_{t_0+h} - M^S_{t_0}|] &= \E_\P[\E_\P[|M^S_{t_0+h} - M^S_{t_0}|  \big | M^S_{t_0}]] = \int_{-\infty}^\infty m_1\biggl(\pi_x^{S ,t_0 ,t_0 +h}\biggr) \, d\mu_{t_0}(x) 
\\ \label{n1/2}	 &=    \int_{-\infty}^\infty F(x) \, d\mu_{t_0}(x)\le  h^{\frac{1}{2}},
\end{align}
where, alleviating the notation from $\pi_x^{S ,t_0 ,t_0 +h}$ to $\pi_x$, and
\[ F(x)=\int_{-\infty}^\infty |x - y| \, d\pi_x(y). \]

\textit{Claim:} The function $x \mapsto F(x)$ satisfies the estimate
\begin{equation}\label{oneLip}
	|F(x) - F(x+k)| \le 2k,  \qquad x \in \R,\quad k>0.
\end{equation}
Indeed,

\begin{align*}
	|F(x+k) - F(x)| &= \left| \int_{-\infty}^{\infty} |y - (x + k)| \, d\pi_{x+k}(y) - \int_{-\infty}^{\infty} |y-x| \, \pi_x(y) \right|
	\\ &\leq  k + \left| \int_{-\infty}^\infty |y - (x+k)| \, d\pi_{x}(y) - \int_{-\infty}^\infty |y-x| \, d\pi_x(y) \right| \leq 2k
\end{align*}
proving the claim. 
In the first inequality we have used the fact that $\pi$ is a Lipschitz Markov kernel.

Choose a compact interval $[a,b]$ such that $K \subseteq [a+1,b-1]$ and denote by $l>0$ a lower bound for the density function $p_{t_0}$ of $\mu_{t_0}$ on $[a,b]$.

We want to estimate $F(x_0)$, for $x_0 \in K$, and start by showing the rough estimate $F(x_0) \le  2$.
Using (\ref{oneLip}) we otherwise have $\int_{x_0}^{x_0+1} F(x) \, dx \ge 1$ and we arrive at the following contradiction to (\ref{n1/2}) for small enough $h$
\[ 1 \le \int_{x_0}^{x_0+1} F(x) \, dx  \le \frac{1}{l} \int_{x_0}^{x_0+1}F(x) \, dp_{t_0}(x) dx \le \frac{1}{l} h^{\frac{1}{2}}. \]



If $F(x_0) \le 2$ we may argue similarly, using again
(\ref{oneLip}) and elementary geometry, to obtain, for $x_0 \in K$,
\[ \frac{1}{8} (F(x_0))^2 \le \int_{x_0}^{x_0+1} F(x) \, dx \le \frac{1}{l} \int_{x_0}^{x_0+1}F(x) \, d\mu_{t_0}(x) \le \frac{1}{l} h^{\frac{1}{2}},\]
yielding the desired estimate for $h$ sufficiently small
\[ \sup_{x\in K}   m_1(\pi_x) =   \sup_{x\in K}  F(x) \le D h^{\frac{1}{4}}, \]
where the constant $D>0$ only depends on the compact set $K$, but not on $h$.

Finally we have to come back to our assumption $I=\R$ which allowed us to imbed the compact set $K$ into the interval $[a,b] \subset I$ such that $ K \subset [a+1,b-1]$. If $I$ is only one- or two-sided bounded we have to reason slightly more carefully, as we can imbed the compact set $K$ only into an interval $[a,b] \subset I$ such that $[a + \epsilon,b- \epsilon]$ contains $K$. But no difficulties arise from replacing $1$ by $\epsilon$ and it is straightforward to adapt the above argument also to this situation.
\end{proof}

% \begin{lemma}\label{lemma1}
% Under the assumptions of Lemma \ref{sublemma} we obtain the following estimate for the square roots of the second moments of
% $\pi_x^{S ,t_0 ,t_0 +h}$.

% \begin{equation} \label{m2}
% m_2\biggl(\pi_x^{S ,t_0 ,t_0 +h}\biggr) := \E_\P\left[\left|M^S_{t_0 +h} - M^S_{t_0}\right|^2  \biggm| M^S_{t_0}=x\right]^{\frac{1}{2}} 
% =\left(\int_{-\infty}^\infty \lvert y-x\rvert^2 \, d\pi_x^{S ,t_0 ,t_0 +h}(y)\right)^{\frac{1}{2}}
%  \le D h^{\frac{1}{4}}, \quad x \in K.
% \end{equation}

% %\begin{equation}\label {m2}
% %m_2(\pi_x) := \E_{\pi_x}[(M_{\frac{k-1}{n}} - M_{\frac{k}{n}})^2] \le C n^{-\frac{1}{2}}, \qquad x \in K.
% %\end {equation}

% \end{lemma}


% %Den (einfachen) Beweis hab ich wo aufgeschrieben, finde ihn aber nicht. Den Exponenten muss ich auch checken.

% \comment{Diesen Beweis habe ich sicherheitshalber aufgehoben:}

% \begin{proof} of Lemma \ref{lemma1}

% First suppose $I=\R$ as in Lemma \ref{sublemma}.

% Again by (\ref{1/n}) we have

% $$\E_{\mu_{\frac{k-1}{n}}}[|M_{\frac{k-1}{n}} - M_{\frac{k}{n}}|^2] = n^{-1}.$$
% We may rewrite this equality in the form 

% $$\E_{\mu_{\frac{k-1}{n}}}[|M_{\frac{k-1}{n}} - M_{\frac{k}{n}}|^2] = \int_{-\infty}^\infty (G_1(x) + G_2(x))d \mu_{\frac{k-1}{n}}(x)=  n^{-1},$$
% where

% $$G_1(x)=\int_x^\infty(y-x)^2d\pi_x(y),$$
% $$G_2(x)=\int^x_{-\infty}(y-x)^2d\pi_x(y).$$


% \textit{Claim:} There is a constant $\tilde C>0$ such that , for $0 \le h\le1$, the function $x \to G_1(x)$ satisfies the one-sided Lipschitz estimate

% \begin{equation}\label {twoLip}
% G_1(x) \le G_1(x+h)+\tilde C h,  \qquad x \in \R, 0 \le h\le1.
% \end{equation}
% Indeed,

% $$G_1(x+h) = \int_{x+h}^\infty (y-(x+h))^2d\pi_{x+h}(y)$$

% $$\ge  \int_{x+h}^\infty (y-(x+h))^2d\pi_{x}(y)$$

% $$=  \int_{x+h}^\infty ((y-x)^2 + 2h(y-x) + h^2)d\pi_{x}(y)$$

% $$\ge \int_{x}^\infty ((y-x)^2 +2h(y-x))d\pi_{x}(y) + 2h^2$$

% $$\ge G_1(x) +2Dh + 2h^2,$$
% where again we have used the fact that $\pi_{x+h}$ dominates $\pi_x$ of first order and where the constant $D$ is given by Lemma \ref{sublemma}.

% The remainder of the proof is analogous as in Lemma \ref{sublemma}.

% Finally we have to come back to the assumption $I=\R$. If $I$ is bounded there is nothing to prove so that only the case of a half-open interval $I$ remains. For this case the above reasoning is easily adapted.



% \end{proof}

Under the assumptions of Lemma \ref{sublemma}, Tschebyscheff's inequality and \eqref{m1} allow us to control the difference between medians and means, since for fixed $0<\delta<\frac{1}{4}$
\begin{equation*}
	\pi_x^{S,t_0,t_0 + h}\left( y \colon |y - x| \geq h^{\delta} \right) \leq  Dh^{\frac{1}{4} - \delta}
		% \left\lvert\text{median}(\pi_x^{S,t_0,t_0 + h}) - \text{mean}(\pi_x^{S,t_0,t_0+h})\right\lvert \leq \frac{Dh^{\frac{1}{4}}}{}
\end{equation*}
for $h > 0$ sufficiently small, $x \in K$ and feasible $S \in \mathcal S$. 
Hence, we have in the setting of Lemma \ref{sublemma} that
\begin{equation}\label{mean med control}
	\left\lvert\text{median}(\pi_x^{S,t_0,t_0 + h}) - \text{mean}(\pi_x^{S,t_0,t_0+h})\right\lvert \leq h^{\delta}.
\end{equation}

\begin{lemma} \label{stop}
Let $0 < \delta < \frac{1}{4}$.
Under the assumptions of the Lemma \ref{sublemma} the same conclusion as in (\ref{m1}) holds true for every $[t_0,t_0+h]$-valued stopping time $\tau$ (by possibly changing the constant $D$ to a different $C$) for $x \in K$
\begin{equation} \label{mtau}
	m_1\biggl(\pi_x^{S ,\tau ,t_0 +h}\biggr) := \E_\P\left[\left|M^S_{t_0 +h} - M^S_{\tau}\right|  \biggm| M^S_{\tau}=x\right] 
	=\int_{-\infty}^\infty \lvert y-x\rvert \, d\pi_x^{S ,\tau ,t_0 +h}(y)
	\le C h^{\delta}.
\end{equation}
\end{lemma}

\begin{proof}
	By Corollary \ref{SLM} we have that $x \mapsto \pi^{S,t_0,t}$ is Lipschitz Markov for all $s \in [t_0,t_0+h]$.
	From this, we can deduce the continuity of the map
	\[
		(x,t) \mapsto \int_{-\infty}^{\infty} |z-y| \, d\pi^{S,t,t_0 + h}(z).
	\]
	Therefore, it suffices to show \eqref{mtau} for deterministic $\tau \equiv t \in [t_0,t_0 + h]$ due to the strong Markov property.
	To this end, let $\tilde K$ be a compact interval in $I$, containing the compact set $K$ in its interior and fix the constant $D$ from Lemma \ref{sublemma}, applied to $\tilde K$.

	To do so, find $x \in I$ such that $y$ equals the median of the measure $\pi_x^{S,t_0,t}$, that is
	\[ \pi^{S,t_0,t}_x(]-\infty, y[) \leq \frac{1}{2} \leq \pi^{S,t_0,t}_x(]-\infty, y]).  \]
	Since $x \mapsto \pi_x^{S,t_0,t}$ is Lipschitz Markov by Corollary \ref{SLM}, such an $x$ exists.
	We may use $x$ to obtain the following estimate
	\begin{multline}\label{+}
		\E_\P\left[\left|M^S_{t_0 +h} - y \right| \biggm| M^S_{\tau}=y\right] 
		\\ \le \E_\P\left[\left(M^S_{t_0 +h} - y\right)_+  \biggm| M^S_{\tau} \ge y, M^S_{t_0} = x\right] + \E_\P\left[\left(M^S_{t_0 +h} - y \right)_-  \biggm| M^S_{\tau} \le y, M^S_{t_0} = x\right]
		\\ \le 2 \E_\P\left[ \left| M^{S,t_0,t_0+h} - y \right| \biggm| M^{S}_{t_0} = x \right] \\ \le 2 \E_\P\left[ \left| M^{S,t_0,t_0+h} - x \right| \biggm| M^{S}_{t_0} = x \right] + 2|x - y|.
	\end{multline}
	Here, we used the first item of Corollary \ref{corr} for the first inequality and $y$ being the median of $\pi^{S,t_0,\tau}_x$ for the second.
	Note that for $h$ sufficiently small we have by \eqref{mean med control} that $x \in \tilde K$. 
	Applying the estimates \eqref{m1} and \eqref{mean med control} to \eqref{+} we find
\[
	\E_\P\left[\left|M^S_{t_0 +h} - y \right| \biggm| M^S_{\tau}=y\right] \le 2D h^\frac{1}{4} + 2h^\delta \leq \left( 2D + 2 \right) h^\delta,
\]
which concludes the proof.
\end{proof}

\begin{proof}[Proof of Theorem \ref{cont}]
	We have to combine Proposition \ref{criterion} and Lemma \ref{stop} with a stopping argument.
	To this end, let $(K_n := [a_n,b_n])_{n \in \mathbb N}$ be an increasing sequence of compact intervals exhausting the interval $I = ]a,b[$, where $a,b \in [-\infty, \infty]$. We let $K_1 = \emptyset$.
	Define the stopping times $\tau^{S,n}$ as the first moment when the continuous martingale $M^S$ leaves $K_n$, write $M^{S,n}$ for the stopped process, and denote the differences $M^{S,n+1} - M^{S,n}$ by $\Delta M^{S,n}$.
	

	For the processes $\Delta M^{S,n}$ the assumptions of Proposition \ref{criterion} still hold true as a consequence of Lemma \ref{stop} and the strong Markov property of $M^S$.
	Consider the process
	
	%\comment{Sorry, aber das gef\"allt mir nicht, das $m^S$ in das Martingal hineinzunehmen. Besser: das $ \infty$-dimensionale Martingal als eine ZV betrachten und die ZV $m$ als zweite ZV dazuf\"ugen. Vielleicht noch ein zoom dazu (auch das Ende des Beweis)}
	
	
	\[
		\tilde M^S_t := (\Delta M_t^{S,n})_{n \in \mathbb N}, \quad t \in [0,1],
	\]
	taking values in $\R^\mathbb N$.
	Moreover, let 
	\[
		m^S := \inf \left\{ k \in \mathbb N \colon \left\lvert \Delta M^{S,k} \right\rvert_\infty = 0 \right\}	
	\]
	be the smallest integer such that the entire trajectory of $(M^S_t(\omega))_{0 \le t \le 1}$ is contained in $K_{m^S}$.
	We know already that for every $n\in\mathbb N$, $M^{S,n}$ (and therefore $\Delta M^{S,n}$) allows for f.d.d.\ convergent subnets along $\mathcal S$ where the limits have by Proposition \ref{criterion} continuous versions.
	We want to argue that similarly the pair $(\tilde M^S,m^S)_{S\in\mathcal S}$ taking values in $\mathcal C[0,1]^\mathbb N \times \mathbb N$ admits a convergent subnet, too.
	For this reason, it is sufficient to show the following.

	\emph{Claim:} For every $\epsilon > 0$ there is $n \in \mathbb N$ such that, for every $ S \in \mathcal S$,
	\begin{equation} \label{cont claim}
		\sup_{S \in \mathcal S} \mathbb P\left( m^S > n \right) < \epsilon.
	\end{equation}
	 Indeed, for each $n\in \mathbb N$ (sufficiently large), there are maximal positive numbers $\alpha_n,\beta_n$ with $\alpha_n + \beta_n \le 1$ such that the probability measure
	\[
		\alpha_n \delta_{a_n} + \beta_n \delta_{b_n} + (1 - \alpha_n - \beta_n) \delta_{\frac{\text{mean}(\mu_1) - \alpha_n a_n - \beta_nb_n}{1 - \alpha_n - \beta_n}}
	\]
	is dominated by $\mu_1$ in convex order.
	Since $\mu_1$ puts no mass to the boundary of $I$, there is for each $\epsilon > 0$ an index $N \in \mathbb N$ such that $\alpha_n + \beta_n < \epsilon$ for all $n \ge N$.
	The law of $M^{S,n}$ is dominated in the convex order by $\mu_1$.
	By maximality of $\alpha_n$ and $\beta_n$ we find uniformly for all $S \in \mathcal S$
	\[
		\mathbb P\left( m^{S} > n \right) = \mathbb P\left( \tau^{S,n} < \infty \right) = \mathbb P\left( M^{S,n} \in \{a_n,b_n\} \right) \leq \alpha_n + \beta_n < \epsilon,
	\]
	which yields the claim \eqref{cont claim}.

	By passing to a subnet, still denoted by $\mathcal S$, we thus obtain that $(\tilde M^S, m^S)_{S \in \mathcal S}$ admits a limit $(\tilde M,m)$ w.r.t.\ convergence of finite dimensional distributions, where
	\[ \tilde M_t = (\Delta M_t^n)_{n \in \mathbb N},\qquad t \in [0,1], \]
	and $(\Delta M_t^{S,n})_{S \in \mathcal S}$ has $\Delta M_t^n$ as its f.d.d.\ limit.
	Due to Proposition \ref{criterion} we may choose a version of $\tilde M$ taking values in $\mathcal C[0,1]^\mathbb N$.
	Consider the following process
	\[
		\hat M_t := \sum_{n = 1}^m \Delta M_t^n,	\qquad t \in [0,1],
	\]
	which has continuous trajectories as $m$ is integer-valued finite.
	Note that f.d.d.\ convergence of $\tilde M^S$ to $\tilde M$ yields f.d.d.\ convergence of
	\[
		M^S = \sum_{n = 1}^{m^S} \Delta M^{S,n} \to \sum_{n = 1}^m \Delta M^n = \hat M.
	\]
	We conclude that $\hat M$ and $M$ coincide in law, and thus the Lipschitz Markov martingale $M$ has a version with continuous paths.
\end{proof}


% \section{Continuity of the martingale solution}

% An important question in the present context is the following: under which conditions on a peacock $(\mu_t)_{0\le t \le 1}$, as defined in section 4 above, is there a 
% strong Markov martingale with \textit{continuous trajectories} having the given marginals. We only focus on the one-dimensional case and mention that the corresponding question for higher dimensions remains wide open. The one-dimensional case, however, is fully understood by now, again by the definitive work of Lowther \cite{Lo08b}.

% %\comment{should we also cite other people here}

% %\comment{WS Wir m\''ussen besprechen, ob wir auf SBM abstellen oder sagen, dass wir nur irgendeine stetige Lipschits Markov Lösung für mu und nu brauchen}

% We do not show Lowther's theorem in full generality, but again we want to isolate a sufficient set of assumptions and present a self-contained proof of the continuity theorem based on the concept of stretched Brownian motion. As the notion of stretched Brownian motion makes good sense in higher dimensions too we hope that this approach might open some perspectives to investigate this terra incognita. In the rest of this section we assume that the reader is familiar with the paper \cite{BaBeHuKa20}.

% \begin{assumption}\label{Ass}

% Let $ (\mu_t)_{0 \le t \le 1}$ be a one-dimensional peacock centered at zero, with densities $p_t(x)$ and finite second moments $m^2_2(\mu_t) = \int_{-\infty}^\infty x^2 d\mu_t(x) = \int_{-\infty}^\infty x^2 p_t(x)dx$ such that the function $ t \to m^2_2(\mu_t)$ is continuous. 

% We assume that there is a -- bounded or unbounded -- open interval $I \subset \R$ supporting each $\mu_t$ and such that, for each compact subset $ K \subset I$, the Lebesgue densities $p_t(x)$ of $\mu_t$ are bounded away from zero, uniformly in $x \in K$ and $t \in [0,1]$.

% %We also assume that, for every strike price $k \in I$, the price of the option  $ x \to  \E_{\mu_t}[(x - k)_+]=\int_k^{\infty}xd\mu_t(x)$ is strictly increasing in $t \in [0,1]$.


% %$0 \le t_0 < t_1 \le 1$, the pair $(\mu_{t_0},\mu_{t_1})$ is irreducible in the sense of \cite{BaBeHuKa20}.
% \end{assumption}



% It will be convenient to suppose (w.l.g.\ via a deterministic time change) that $ t \to m^2_2(\mu_t)$ is affine. More precisely we may assume that $m^2_2(\mu_{t+h}) - m^2_2(\mu_t)= h$ so that for all martingales $M$ with $law( M_t) = \mu_t$ and $law( M_{t+h} )= \mu_{t+h}$

% \begin{equation} \label{1/n}
% \E \left[\left(M_{t+h} - M_t \right)^2\right] = h.
% \end{equation}


% \begin{theorem}\label{cont}
% Under the above assumptions there is a \underline{continuous} strong Markov martingale $M = (M_t)_{0 \le t \le 1}$ with given marginals $ (\mu_t)_{0 \le t \le 1}$.
% \end{theorem}

% {\comment {Rot sind die Sachen, die wir jetzt nicht mehr brauchen}}
% {\color {red}We note that $M$ is the \textit{unique} martingale with the above properties by Theorem \ref{uniqueness theorem}. Uniqueness holds actually true in the larger class of \textit{almost continuous} strong Markov martingales (see \cite{Lo08b}).

% %\begin{lemma}\label{lemmab}
% %Assume that, for every compact set $K \subset I$ and $ b > 0$ there is $c>0$ such that every interval $[x,x+b]$ of length $b$ and contained in $K$, has $\mu_t$ measure bigger than $c$, for every $t \in [0,1]$.
 
% %Then the supports of the measures $\mu_t$ are convex, as required by Assumption \ref{Ass}.

% %\end{lemma}

% %\begin {proof}

% %Assume that there is a compact set $K$ in $I$ and $b>0$, as well as a sequence $(x_j)_{j=1}^\infty$ and $(t_j)_{j=1}^\infty$ such that $[x_j, x_j+b]$ is contained in $K$ and $\mu_{t_j} [x_j,x_j+b] < \frac{1}{j}$. We may suppose that $(x_j)_{j=1}^\infty$ converges to some $x_0 \in K$ and $(t_j)_{j=1}^\infty$ to some $t_0 \in [0,1]$. Using the weak convergence of $(\mu_{t_j})_{j=1}^\infty$ to $\mu_{t_0}$ we have $\mu_{t_0} (]x_0,x_0+b[) =0$ so that $\mu_{t_0}$ fails to have convex support.


% %\end {proof}

% %Die Beweise der beiden Lemmas sind noch auf diese neue Annahme anzupassen.

% %Der Beginn des Beweis, den ich zusammengestoppelt habe, muss noch umgeschrieben werden.

% As already mentioned, our approach is to follow the line of argument in \cite[Section 1.4.6]{BaBeHuKa20}. Using the notation of this paper we find, for every $n \in \mathbb N$ and $k \in \{1,\dots,n\}$, a unique martingale $M^{n,k} =(M^{n,k}_t)_{\frac {k-1}{n} \le t \le \frac{k}{n}}$, called stretched Brownian motion from $\mu_\frac{k-1}{n}$ to $\mu_\frac{k}{n}$, such that 
% \[	law\left(M^{n,k}_{\frac{k-1}{n}}\right) = \mu_{\frac{k-1}{n}} \text{ and } law\left(M^{n,k}_{\frac{k}{n}}\right) = \mu_{\frac{k}{n}}.\]
% Clearly these pieces can be pasted together to yield a martingale $M^n = (M^n_t)_{0 \le t \le 1}$ with given marginals $\mu_{\frac{k}{n}}$ at times $\frac{k}{n}$.

% The assumption of strictly increasing option prices is sufficient to guarantee that each  $M^{n,k} =(M^{n,k}_t)_{\frac {k-1}{n} \le t \le \frac{k}{n}}$ is a \textit{standard} stretched Brownian motion. We describe briefly this pleasant structure and refer to \cite{BaBeHuKa20} for details.

% For each fixed $n$ and $k$ there is a probability measure $\alpha^n_\frac{k-1}{n}$ on $\R$  and an increasing  map 
% \[
% 	f^n_{\frac{k}{n}} \colon \R \to \R \text{ such that }\left(f^n_{\frac{k}{n}}\right)_\# \left(\alpha^n_{\frac{k-1}{n}} * \gamma\right) = \mu_{\frac {k}{n}}.
% \]
% Here $\gamma$ denotes a standard centered Gaussian distribution with variance $1$, the symbol $*$ the operation of convolution, and $f_{\#}(\cdot)$ the image of a measure under the map $f$. 
% For $t \in [\frac{k-1}{n},\frac{k}{n}[$ we denote by $f^{n,k}_t$ the map
% \[
% 	f^{n,k}_t(z) := \text{barycenter} \left(\left(f^n_{\frac{k}{n}}\right)_\# (\gamma^z_{k-nt}) \right),
% \]
% where $ \gamma^z_{k-nt}$ denotes the Gaussian with variance $k-nt$ centered at $z$.
% The functions $f^{n,k}_t$ are increasing and we have pointwise $\lim_{t \to \frac{k}{n}} f^{n,k}_{t}  = f^n_{\frac{k}{n}}$. 
% We denote by $\pi^{n,t}_x$ the probablility measure
% \begin{equation}
% 	\label{transition kernel}
% 	\pi^{n,t}_x := \left(f^n_{\frac{k}{n}}\right)_{\#} (\gamma^z_{k-nt}),\qquad \text{where } f^{n,k}_t(z) = x.
% \end{equation}
% In particular, we have the decomposition $\mu_{\frac{k}{n}} = \int_{-\infty}^{\infty}\pi^{n,\frac{k-1}{n}}_x d\mu_{\frac{k-1}{n}}(x)$.

% The crucial issue is that the martingale $M^{n,k}$ can be expressed as the image of the Brownian motion $(B^{n,k}_t)_{\frac{k-1}{n} \le t \le \frac{k}{n}}$, starting with $law(B^{n,k}_{\frac{k-1}{n}} )= \alpha^n_\frac{k-1}{n}$ under the map $(f^{n,k}_t)_{\frac{k-1}{n} \le t \le \frac{k}{n}}$ via 
% \begin{equation}\label{Brown}
% 	M^{n,k}_t = f^{n,k}_t(B^{n,k}_t), \quad \frac{k-1}{n} \le t < \frac{k}{n}, \qquad \text{and }M_\frac{k}{n}^{n,k} = f_\frac{k}{n}^n(B_\frac{k}{n}^{n,k}).
% \end{equation}
% %  and $M^{n,k}_{\frac{k}{n}} =  f^n_{\frac{k}{n}}(B^{n,k}_\frac{k}{n})$ (see \cite{BaBeHuKa20}). 
%  The probability measure $\alpha^n_{\frac{k-1}{n}}$ on $\R$ is uniquely determined, up to a parallel shift by a constant. 
% %\comment{Should we leave out these two sentences?}
% %We also mention that in view of ($\ref{1/n}$) it is natural in our context to use the normalized Gaussian $\gamma_{\frac{1}{n}}$  with variance $\frac{1}{n}$ while in [BBHL] the Gaussian $\gamma$ with variance $1$ is considered. The only difference caused by this normalization is a stretching of the measure $\alpha^n_\frac{k-1}{n}$ and the function $f^n_{\frac{k-1}{n}}$ by a factor of $n^{\frac{1}{2}}$.


% We thus have that each $M^n = (M^n_t)_{0 \le t \le 1}$ is a continuous strong Markov martingale and therefore Lipschitz Markov (Corollary \ref{SLM}). }



% There is an obvious and well-known strategy for the proof. We want to obtain the desired martingale $M$ as a limit of approximations which fit the peacock $(\mu_t)_{0 \le t \le 1}$ on finitely many points of time. As in \cite{HiRoYo14} it is convenient to do so along the ordered set $\cal S$ of finite partitions $ S = \{s_0,\dots,s_n\} $ of $[0,1]$, where $0 = s_0 < \dots < s_n =1$.

% For each $S \in \cal S$ we choose a continuos strong Markov martingale $M^S$ having the given marginals at each time $s_i \in S$. This is standard by now and there are a number of suitable constructions in the literature XXXXX  For example SBM XXXXXX{\comment {Es schaut so aus, dass wir f\''ur den Beweis SBM nicht brauchen :-((}}

% Identifying the martingales $M^S$ with their induced measures on path space  $\mathcal C[0,1]$  this family of measures clearly is tight, if considered on Skorohod space $\mathcal D[0,1]$, equipped with the topology of convergence of finite dimensional distributions. Hence we can find a cluster point $M$ in the set $\mathcal P(\mathcal D[0,1])$ of probability measures on $\mathcal D[0,1]$, see e.g.\ \cite{BeHuSt16} for the straightforward argument. By refining the filter $\mathcal S$ we may suppose that $M$ is a limit point. 

% We fix such a limiting process $M$ which by  Proposition \ref{LMclosed} is a Lipschitz Markov martingale.% It was the path-breaking insight of Kellerer \cite{Ke72, Ke73} that the Lipschitz Markov property allows to conclude that $M$ again is a strong Markov martingale (still being Lipschitz Markov). 
%  These arguments are standard by now and, e.g., well presented in the papers \cite{HiRoYo14, BeHuSt16}.  A priori the martingale $M$ has c\`adl\`ag trajectories.
% Our present task is to show the \textit{continuity} of the trajectories of the limiting process $M$ under the above assumptions.

% We first give a general criterion for the continuity of a limiting martingale $M$ which is somewhat reminiscent of the classical Kolmogorov continuity criterion XXXXX. 
% %We switch from the notation $M^S$ to $M^\alpha$ to stress te generality of the argument.

% \begin{proposition}\label{criterion}

% Let $(M^\alpha)_{\alpha \in I}$ be a net of $\R$-valued continuous strong Markov martingales  $ M^\alpha = (M^\alpha_t)_{0 \le t \le 1}$ and $M$ its limit in $\mathcal P(\mathcal D[0,1])$ with respect to convergence in finite dimensional distribution.

% Suppose that there are constants $C_1>0$ and $\beta > 0$ such that, for every $0 \le t_0 < t_0 + h \le 1$ and every $\alpha \in I$

% \begin {equation}\label{BMO}
% ||(M^\alpha_t)_{t_0 \le t \le t_0 + h}||_{\textit{BMO}_1} := \sup_\tau\{ || \E[|M^\alpha_{t_0 + h} -M^\alpha_\tau| \Big{|} {\cal{F}}_{\tau}] ||_\infty   \}   \le  C_1 h^\beta.
% \end{equation}
% Then the martingale $M$ has continuous trajectories.

% \end{proposition}

% In (\ref{BMO}) $\tau$ runs through the $[t_0,t_0 + h]$-valued stopping times with respect to the natural filtration of $ (M^\alpha_t)_{t_0 \le t \le t_0 + h}$. As $M^\alpha$ is strong Markov, condition (\ref{BMO}) is tantamount to the requirement that the first moments $m_1(\pi_x^{\alpha ,\tau ,t_0 +h}) = \int_{-\infty}^\infty |y-x| d\pi_x^{\alpha ,\tau ,t_0 +h}(y)$ of the transition probabilities $\pi_x^{\alpha,\tau,t_0+h}$ from $M^\alpha_\tau = x$ to $M^\alpha_{t_0 + h}$ satisfy 

% \begin{equation}\label {beta}
%  m_1( \pi_x^{\alpha,\tau,t_0 +h})  = \int_{-\infty}^\infty |y-x| d\pi_x^{\alpha ,\tau ,t_0 +h}(y) \le C_1 h^\beta, 
%  \end{equation}
%  for $\mu^\alpha_\tau$-almost all $x \in \R$, where $\mu^\alpha_\tau$ denotes the law of $M^\alpha _\tau$.



% An important feature of the $BMO$ norms for \textit{continuous} martingales is that, by  the John-Nirenberg inequality, all $BMO_q$ norms are equivalent, for $1 \le q < \infty$ (see, e.\ g.\ [Kazamaki]). Applying this fact to the present context, inequality (\ref{BMO}) is equivalent to the existence of a constant $C_q > 0$ (for some or, equivalently, for all $1 \le q < \infty$) such that

% \begin {equation}\label{BMOq}
% ||(M^\alpha_t)_{t_0 \le t \le t_0 + h}||_{\textit{BMO}_q} := \sup_\tau\{ || \E[|M^\alpha_{t_0 + h} -M^\alpha_\tau|^q \Big{|} {\cal{F}}_{\tau}]^{\frac{1}{q}} ||_\infty   \}   \le  C_q h^\beta.
% \end{equation}
% Again this can be rephrased as 

% \begin{equation}\label {betaq}
%  m_q(\pi_x^{\alpha ,t_0 ,t_0 +h}) = \left(\int_{-\infty}^\infty |y-x|^q d \pi_x^{\alpha ,t_0 ,t_0 +h}(y)\right)^{\frac{1}{q}} \le C_q h^\beta. 
%  \end{equation}

% \begin {proof} of Proposition \ref{criterion}.

% Suppose that $M$ fails to be continuous and let us work towards a contradiction. If $M$ has jumps of size bigger than $2a > 0$ with probability bigger than $\kappa >0$ we can find, for sufficiently small $h >0$, times $ 0 \le t_0 < t_0 +h \le 1$ such that 

% \begin {equation}\label {a}
% \P[|M_{t_0 +h} - M_{t_0}| > a ] > h\kappa.
% \end {equation}
% The same inequality holds true for $M$ replaced by $M^\alpha$, for large enough $\alpha$, in view of the convergence in finite dimensional distribution. Fixing such an $\alpha$, it follows that there is a set $A \subset \R$ of positive measure with respect to the law of $M^\alpha_{t_0}$ such that, for $x \in A$,

% \begin {equation}\label {acond}
% \P[|M^\alpha_{t_0 +h} - M^\alpha_{t_0}| > a  \big| M^\alpha_{t_0} = x] > h\kappa.
% \end {equation}
% For $x \in A$ we therefore have 

% \begin {equation}\label {mq}
% m_q(\pi_x^{\alpha ,t_0 ,t_0 +h}) = \left(\int_{-\infty}^\infty |y-x|^q d \pi_x^{\alpha ,t_0 ,t_0 +h}(y)\right)^{\frac{1}{q}} > (ah\kappa )^{\frac{1}{q}}.
% \end {equation}
% Choosing $1\le q < \infty$ such that $q > \frac{1}{\beta}$, we arrive at the desired contradiction to (\ref{BMOq}), where we let $\tau = t_0$.     \end {proof}


% Turning back to the above defined family $(M^S)_{S \in \cal S}$ of martingales we shall establish an inequality of the type (\ref{beta}) for the transition probabilities $\pi_x^{S ,t_0 ,t_0 +h}$, using the fact that the functions $x \mapsto \pi_x^{S ,t_0 ,t_0 +h} $ are non-decreasing with respect to first order dominance.



% %\begin{lemma}\label{lemma1}
% %For a compact set $K \subset I$ there is a constant $C>0$ such that, for each $n \in \mathbb N$ and $k = 1,\dots,n$, the second moments of the transition measures $\pi_x$ can be estimated by

% %\begin{equation}\label {m2}
% %m_2(\pi_x) := \E_{\pi_x}[(M_{\frac{k-1}{n}} - M_{\frac{k}{n}})^2] \le C n^{-\frac{1}{2}}, \qquad x \in K.
% %\end {equation}

% %\end{lemma}

% %We first establish an analogous  preparatory result.

% \begin{lemma}\label{sublemma}
% Let $(\mu_t)_{0 \le t \le 1}$ be a peacock satisfying Assumption \ref{Ass}. Fix $0 \le t_0 < t_0+h \le 1$ and consider a martingale $M^S$ as above such that $S = \{s_0, \dots, s_n\}$ contains $t_0$ and $t_0 + h$.

% For a compact set $K \subset I$ there is a constant $D>0$ such that the first moments of the transition measures $\pi_x^{S ,t_0 ,t_0 +h}$ can be estimated by

% \begin{equation} \label{m1}
% m_1\biggl(\pi_x^{S ,t_0 ,t_0 +h}\biggr) := \E_\P\left[\left|M^S_{t_0 +h} - M^S_{t_0}\right|  \biggm| M^S_{t_0}=x\right] 
% =\int_{-\infty}^\infty \lvert y-x\rvert \, d\pi_x^{S ,t_0 ,t_0 +h}(y)
%  \le D h^{\frac{1}{4}}, \quad x \in K,
% \end{equation}
% where $\P$ denotes the law of the martingale $M^S$.

% \end{lemma}

% \begin{proof}
% We first suppose that $I=\R$. 

% By (\ref{1/n}) we have
% \[
% 	\E_\P[|M^S_{t_0+h} - M^S_{t_0}|] \le  h^{\frac{1}{2}}.
% \]
%  We may rewrite this inequality in the form 
% \begin{multline}\label{n1/2}
% \E_\P[|M^S_{t_0+h} - M^S_{t_0}|] = \E_\P[\E_\P[|M^S_{t_0+h} - M^S_{t_0}|  \big | M^S_{t_0}]] =
% \\	\int_{-\infty}^\infty m_1\biggl(\pi_x^{S ,t_0 ,t_0 +h}\biggr) \, d\mu_{t_0}(x) =    \int_{-\infty}^\infty \left (F_1(x) + F_2(x) \right )  \, d\mu_{t_0}(x)\le  h^{\frac{1}{2}},
% \end{multline}
% where, alleviating the notation from $\pi_x^{S ,t_0 ,t_0 +h}$ to $\pi_x$, 
% \[ F_1(x)=\int_x^\infty(y-x) \, d\pi_x(y) \text{ and } F_2(x)=\int^x_{-\infty}(x-y) \, d\pi_x(y). \]

% \textit{Claim:} The function $x \to F_1(x)$ satisfies the one-sided Lipschitz estimate
% \begin{equation}\label{oneLip}
% 	F_1(x) \le F_1(x+k)+k,  \qquad x \in \R,\quad k>0.
% \end{equation}
% Indeed,
% \begin{align*}
% 	F_1(x+k) &= \int_{x+k}^\infty y-(x+k) \, d\pi_{x+k}(y) \ge  \int_{x+k}^\infty y-(x+k) \, d\pi_{x}(y)
% \\	&\ge\int_{x}^\infty y-(x+k) \, d\pi_{x}(y) \ge \int_{x}^\infty (y-x)d\pi_{x}(y) -k,
% \end{align*}
% proving the claim. In the first inequality we have used the fact that $\pi_{x+k}$ dominates $\pi_x$ of first order.

% Choose a compact interval $[a,b]$ such that $K \subseteq [a+1,b-1]$ and denote by $l>0$ a lower bound for the density function $p_{t_0}$ of $\mu_{t_0}$ on $[a,b]$.

% We want to estimate $F_1(x_0)$, for $x_0 \in K$, and start by showing the rough estimate $F_1(x_0) \le  1$. Using (\ref{oneLip}) we otherwise have $\int_{x_0}^{x_0+1} F_1(x) dx \ge \frac{1}{2}$ and we arrive at the following contradiction to (\ref{n1/2}) for small enough $h$

% \[ \frac{1}{2} \le \int_{x_0}^{x_0+1} F_1(x) \, dx  \le \frac{1}{l} \int_{x_0}^{x_0+1}F_1(x) \, dp_{t_0}(x) dx \le \frac{1}{l} h^{\frac{1}{2}}. \]



% If $F_1(x_0) \le 1$ we may argue similarly, using again
% (\ref{oneLip}) and elementary geometry, to obtain, for $x_0 \in K$,
% \[ \frac{1}{2} (F_1(x_0))^2 \le \int_{x_0}^{x_0+1} F_1(x) \, dx \le \frac{1}{l} \int_{x_0}^{x_0+1}F_1(x) \, d\mu_{t_0}(x) \le \frac{1}{l} h^{\frac{1}{2}}. \]
% The same reasoning applies to $F_2$ yielding the desired estimate
% \[ \sup_{x\in K}   m_1(\pi_x) =   \sup_{x\in K} \{F_1(x) + F_2(x)\} \le D h^{\frac{1}{4}}, \]
% for a constant $D>0$ which only depends on the compact set $K$, but not on $h$.

% Finally we have to come back to our assumption $I=\R$ which allowed us to imbed the compact set $K$ into the interval $[a,b] \subset I$ such that $ K \subset [a+1,b-1]$. If $I$ is only one- or two-sided bounded we have to reason slightly more carefully, as we can imbed the compact set $K$ only into an interval $[a,b] \subset I$ such that $[a + \epsilon,b- \epsilon]$ contains $K$. But no difficulties arise from replacing $1$ by $\epsilon$ and it is straightforward to adapt the above argument also to this situation.
% \end{proof}

% \begin{lemma}\label{lemma1}
% Under the assumptions of Lemma \ref{sublemma} we obtain the following estimate for the square roots of the second moments of
% $\pi_x^{S ,t_0 ,t_0 +h}$.

% \begin{equation} \label{m2}
% m_2\biggl(\pi_x^{S ,t_0 ,t_0 +h}\biggr) := \E_\P\left[\left|M^S_{t_0 +h} - M^S_{t_0}\right|^2  \biggm| M^S_{t_0}=x\right]^{\frac{1}{2}} 
% =\left(\int_{-\infty}^\infty \lvert y-x\rvert^2 \, d\pi_x^{S ,t_0 ,t_0 +h}(y)\right)^{\frac{1}{2}}
%  \le D h^{\frac{1}{4}}, \quad x \in K.
% \end{equation}

% %\begin{equation}\label {m2}
% %m_2(\pi_x) := \E_{\pi_x}[(M_{\frac{k-1}{n}} - M_{\frac{k}{n}})^2] \le C n^{-\frac{1}{2}}, \qquad x \in K.
% %\end {equation}

% \end{lemma}

% \begin{proof}   Den (einfachen) Beweis hab ich wo aufgeschrieben, finde ihn aber nicht. Den Exponenten muss ich auch checken.

% \end{proof}

% Lemma \ref{lemma1} and the elementary inequality $|\textit{barycenter}(\pi) - \textit{median}(\pi)| \le m_2( \pi)$ will allow us to estimate the difference between  the medians and the means appearing in the proof of the subsequent lemma.

% \begin{lemma} \label{stop}

% Under the assumptions of the Lemma \ref{sublemma} the same conclusion as in (\ref{m1}) holds true for every $[t_0,t_0+h]$-valued stopping time $\tau$ (by possibly changing the constant $D$ to a different $C$).

% \begin{equation} \label{mtau}
% m_1\biggl(\pi_x^{S ,\tau ,t_0 +h}\biggr) := \E_\P\left[\left|M^S_{t_0 +h} - M^S_{\tau}\right|  \biggm| M^S_{\tau}=x\right] 
% =\int_{-\infty}^\infty \lvert y-x\rvert \, d\pi_x^{S ,\tau ,t_0 +h}(y)
%  \le C h^{\frac{1}{4}}, \quad x \in K,
% \end{equation}


% \end{lemma}

% \begin{proof}

% Let $\tilde K$ be a compact interval in $I$, containing the compact set $K$ in its interior and fix the constant $D$ from Lemma \ref{sublemma}, applied to $\tilde K$.

% Suppose first that the stopping time $\tau \in [t_0,t_0+h]$ is deterministic. To estimate $m_1\biggl(\pi_y^{S ,\tau ,t_0 +h}\biggr)$ it will suffice -- up to a factor of $2$ -- to estimate $\E_\P\left[\left(M^S_{t_0 +h} - M^S_{\tau}\right)_+  \biggm| M^S_{\tau}=y\right] $, similarly as in the proof of Lemma \ref{sublemma}.

% To do so, find $x \in \R$ such that $y$ equals the median of the measure $\pi_x^{S,t_0,\tau}$. As $M^S$ is a Feller process the function $ x \mapsto \pi_x^{S,t_0,\tau}$ is weakly continuous so that  such an $x$ exists. To be somewhat formalistic: in general, the median of $\pi_x^{S,t_0,\tau}$ is defined as a non-empty interval and we can choose any $x$ in this interval. From Lemma \ref{stop} we obtain the estimate

% \begin{equation}\label{xy}
% |x-y| \le Dh^{\frac{1}{4}}.
% \end{equation}

% Using Corollary \ref{corr} (i) we obtain the crucial estimate


% \begin{equation}\label{+}
% \E_\P\left[\left(M^S_{t_0 +h} - y \right)_+  \biggm| M^S_{\tau}=y\right] \le 2 \E_\P\left[\left(M^S_{t_0 +h} - y \right)_+  \biggm| M^S_{\tau} \ge y, M^S_{t_0} = x\right] \le
% \end{equation}

% $$ 2 \E_\P\left[\left(M^S_{t_0 +h} - y\right)_+  \biggm| M^S_{t_0} = x\right] 
%  $$
 
%  Indeed, conditioning on $M^S_{t_0} =x$, we know that $\P [M^S_\tau \ge y | M^S_{t_0} =0] \ge \frac{1}{2}$, as $y$ is the mean of $\pi_x^{S,t_0,\tau}$. By Corollary \ref{corr} (i) and noting that the function $ z \mapsto (z-y)_+ $ is non-negative we obtain (\ref{+}).
 
%  %We still have to take care of the difference of the mean $x$ and the median $y$ of the measure $\pi_x^{S,t_0,\tau}$. This is provided by Lemma \ref{lemma1} and the above mentioned mean-median inequality. 
%  Using(\ref{xy}) we may continue to estimate
 
 
%  $$2 \E_\P\left[\left(M^S_{t_0 +h} - y\right)_+  \biggm| M^S_{t_0} = x\right] 
%  \le 2 \E_\P\left[\left(M^S_{t_0 +h} - x\right)_+  \biggm|M^S_{t_0} = x\right] + 2|x-y| \le $$
 
%  $$m_1( \pi^{S,t_0,t_0+h}_x) +2D h^{\frac{1}{4}}  \le 3D h^{\frac{1}{4}}. $$


% In conclusion, we obtain (\ref{mtau}) for deterministic stopping times $\tau \in [t_0,t_0+h]$ for some uniform constant $C$. The same inequality holds true for stopping times taking finitely many values in $[t_0,t_0+h]$. Using once again the Feller property of $M^S$ we also get the conclusion for general $[t_0,t_0+h]$-valued stopping times $\tau$.

% $$ \le 2 \E_\P\left[\left(M^S_{t_0 +h} - x\right)_+  \biggm|M^S_{t_0} = x\right] + |x-y| \le m_1( \pi^{S,t_0,t_0+h}_x) +|x-y|. $$




% \end{proof}


% {\color {red}   \begin{lemma}\label{dist}
% Fix a compact convex set $ K\subset I$ and $ 0 < \epsilon < \frac{1}{4}$.
% For $n$ large enough, $k = 1,\dots,n$, and $x_1,x_2 \in K$, let $z_1,z_2 \in \R$ be such that $f^n_{\frac{k}{n}} (z_1) = x_1$ and $f^n_{\frac{k}{n}} (z_2) = x_2$. If $|x_1 -x_2| \geq n^{-\frac{1}{4} + \epsilon}$ we have
% \[ |z_1 -z_2| \ge 1.\]
% Hence, for all $x_1,x_2 \in K$ with $|x_1 -x_2| \ge a > 0$, we have $|z_1 - z_2| \ge \lfloor an^{\frac{1}{4} - \epsilon}\rfloor$.
% \end{lemma}

% \begin{proof}
% Applying Tschebyscheff's inequality to Lemma \ref{m1} we have for $0 < \delta < \epsilon$
% \[ \lim_{n \to \infty}\pi^{n,\frac{k-1}{n}}_x\left( \left\{ y \in \R : |y-x| \ge n^{-\frac{1}{4} + \delta} \right\} \right) = 0, \]
% uniformly in $x\in K$ and $k = 1,\dots,n$.

% If $|x_1 -x_2| \geq n^{-\frac{1}{4} + \epsilon} > 2n^{-\frac{1}{4} + \delta}$ we can bound the total mass of the minimum of the two probability measures $\pi^{n,\frac{k-1}{n}_{x_1}}$ and $\pi^{n,\frac{k-1}{n}}_{x_2}$ by
% \begin{multline*}
% 	\pi^{n,\frac{k-1}{n}}_{x_1} \wedge  \pi^{n,\frac{k-1}{n}}_{x_2}(\R) \leq 
% \\	\pi^{n,\frac{k-1}{n}}_{x_1}\left( \left\{ y \in \R : |y-x_1| \ge n^{-\frac{1}{4} + \delta} \right\}  \right) \wedge  \pi^{n,\frac{k-1}{n}}_{x_2}\left( \left\{ y \in \R : |y-x_1| \ge n^{-\frac{1}{4} + \delta} \right\}  \right)
% \\	+	\pi^{n,\frac{k-1}{n}}_{x_1}\left( \left\{ y \in \R : |y-x_2| \ge n^{-\frac{1}{4} + \delta} \right\}  \right) \wedge  \pi^{n,\frac{k-1}{n}}_{x_2}\left( \left\{ y \in \R : |y-x_2| \ge n^{-\frac{1}{4} + \delta} \right\}  \right),
% \end{multline*}
% and therefore is smaller than, say, $1/100$, for $n$ large enough.
% By \eqref{transition kernel} we have that the total mass of the measure $\gamma^{z_1} \wedge \gamma^{z_2}$ equals the total mass of $\pi^{n,\frac{k-1}{n}}_{x_1}   \wedge  \pi^{n,\frac{k-1}{n}}_{x_2}$, and is consequentially less than $1/100$.
% In conclusion, the distance of the centers $z_1$ and $z_2$ of the two Gaussian distributions $\gamma^{z_1}$ and $\gamma^{z_2}$ is bigger than $1$.

% The final statement of the lemma follows from chopping the interval $[x_1,x_2]$ into at least $\lfloor a n^{\frac{1}{4} - \epsilon} \rfloor$ pieces of length greater than $n^{-\frac{1}{4} + \epsilon}$ and making use of monotonicity of the maps $f^n_{\frac{k}{n}}$.
% \end{proof}

% %\begin{remark}

% %\end{remark}\label{rem}
% %Let $t \in [\frac{k-1}{n}, \frac{k}{n}]$ and denote by $\pi^t_x$ the law of $M^n_{\frac{k}{n}}$, conditionally on $M^n_t = x$. An inspection of the proof reveals that the conclusions of Lemma \ref{m2} (as well as of Lemma \ref{m1}) still hold true with the same constants if we replace $\pi_x$ by $\pi^t_x$.


% \bigskip   }

% \begin{proof}[Proof of Theorem \ref{cont}]

% {\color{red}

% Assume that the limiting martingale $M$ fails to be continuous and let us work towards a contradiction to Assumption $\ref {Ass}$. 
% Suppose that the trajectories of the limiting process $M$ have jumps (w.l.g.\ upward jumps) of size at least $2a > 0$  with strictly positive probability,
% that means
% \[
% 	\mathbb P\left( \left\{ \text{The trajectory of } M \text{ has an upward jump of size greater }2a \right\} \right) \geq 2 \kappa > 0,
% \]
% for $\kappa > 0$.
% Let $\tau$ be the time of the first upward jump of size at least $2a$
% \[ \tau(\omega) := \inf \left\{ t > 0 \colon M_t(\omega) - \lim_{s\nearrow t} M_s(\omega) \geq 2a\right\}. \]


% We then may find a convex and compact subset $K \subset I$ and $\kappa>0$ such that, for each $n \in \mathbb N$ there is some $k \in \{1,\dots,n\}$ such that

% \begin{equation}\label{jump1} 
% 	\P[M^{n,k}_{\frac{k}{n}}  \ge M^{n,k}_{\frac{k-1}{n}} +a, M^{n,k}_{\frac{k-1}{n}} \in K, M^{n,k}_{\frac{k}{n}} \in K] \ge \frac{\kappa}{n}.
% \end{equation}
% We thus may find $x \in K$ such that $x+a$ is in $K$ too, and such that

% \begin{equation}\label{jump}
% \P[M^{n,k}_{\frac{k}{n}}  \ge x + a \quad | \quad M^{n,k}_{\frac{k-1}{n}} = x] \ge \frac{\kappa}{n}.
% \end{equation}

% Define $z$ and $z'$ via  $f^n_{\frac{k}{n}}(z) = x$ and $f^n_{\frac{k}{n}}(z') = x+a$. From Lemma \ref{dist} we deduce that $|z - z'| > a n ^{\frac{1}{4} - \epsilon}$, for $0<\epsilon<\frac{1}{4}$ and $n$ big enough.



% Consider the martingale $(M_t^{n,k,x})_{\frac{k-1}{n} \le t \le \frac{k}{n}}$ conditioned to start at $M^{n,k,x}_{\frac{k-1}{n}} = x$. By (\ref{jump}) we know that $\P[M^{n,k,x}_{\frac{k}{n}}  \ge x + a ] \ge \frac{\kappa}{n}$. This translates to the behavior of the Brownian motion $B^{n,k,z}$ via (\ref{Brown}), conditioned to start at $ B^{n,k,z}_{\frac{k-1}{n}} = z$. We thus obtain $\P[B^{n,k,z}_{\frac{k}{n}}  \ge z' ] \ge \frac{\kappa}{n}$.

% We thus arrive in view of the last statement of Lemma \ref{dist} at the desired contradiction: it is clearly absurd that, for $\epsilon < \frac{1}{4}$, a standard Brownian motion travels as far as $n ^{-\frac{1}{4} - \epsilon}$ with probability $\frac{\kappa}{n}$ during an interval of time of length $1$.      }

% We only have to combine Proposition \ref{criterion} and Lemma \ref{stop} with a simple stopping argument. Fix a compact interval $K \subset I$, define the stopping times $\tau^{S,K}$ as the first moment when the continuous martingale $M^S$ leaves $K$, and consider the processes $M^{S,\tau^{S,K}}$, stopped at time $\tau^{S,K}$. 

% For the stopped processes $M^{\tau^{S,K}}$ Proposition \ref{criterion} still holds true. Again we may find a  limit $M^K$ along the ordered set $\mathcal S$ (possibly refining it, if necessary), and by Lemma \ref{stop} this limit $M^K$ is a continuous strong Markov martingale. Clearly we want to show that $M^K$ converges to $M$ in an appropriate sense which suffices to make sure that $M$ is continuous.

% To do so, note that the distribution of the terminal value $M_1^S$ is identical for all $S \in \mathcal S$, namely $\mu_1$. Hence, for $\epsilon >0$ , we may find a large enough compact interval $K \subset I$ such that $\P [\tau^{S,K} = \infty] < \epsilon$, for all $S \in \mathcal S$. It follows that $||M^S_1 - M^{S,K}_1||_{L^2(\P)}$ tends to zero, uniformly in $S \in \mathcal S$. Hence $||M_1 - M^{K}_1||_{L^2(\P)}$ also tends to zero so that, by Doob's maximal $L^2(\P)$ inequality, we may indeed conclude that the limiting process $M$ is continuous.




% %Choosing a sequence $(K_n)_{n =1}^\infty \subset I$ so that $\mu_1[K_n]$ tends to $1$, we obtain from Doob's $L^2$ maximal inequality that $(M^{K_n})_{n=1}^\infty$ converges with respect to the 


% \end{proof}
% %Let $(t_j)_{j=0}^J$ be increasing in $[x,x+2a]$ such that the distances  satisfy $t_j - t_{j-1} \ge n^{-\frac{1}{4}}$. We can choose these points such that $J >  n^{\frac{1}{4}}a$.

% %Define the stopping times $(\tau_j)_{k=0}^J$ as the first hitting of the level $t_j$ by the martingale $(M_t^{n,x})_{\frac{k-1}{n} \le t \le \frac{j}{n}}$ where, as usual, we set $\tau_j = \infty$ if this does not happen.

% %The crucial observation is that by Lemma \ref{m2} and Remark \ref{rem} we have for $n$ big enough 

% %\begin{equation}\label {Doob}
%  %\P  \{ \tau_j < \infty | \tau_{j-1} < \infty \} < \frac{1}{2}, \qquad j = 1,\dots,J.   
%  %\end{equation}
% %Indeed, the second moment of $\pi^{\tau_{j-1}}_{t_{j-1}}$ is of the order of $ n^{-1}$. As the distance from $t_{j-1}$ to $t_j$ is of the order $n^{-\frac{1}{4}}$, we obtain from Doob's maximal inequality for the $L^2$ norm and Chebysheff's inequality that the probability that $(M^n_t)_{\tau_{j-1} \le t \le \frac{k}{n}}$ moves from $t_{j-1}$ to $t_j$ tends to zero, uniformly in $ k = 1,\dots,n$, as $n$ goes to infinity.

% %Intern: wir k\''onnen hier auch die schwache $L^1$ Doob Ungleichung verwenden und Lemma \ref{m1}. Aber die $L^2$ norm ist aus anderen Gr\''unden nat\''urlich im Beweis.

% %Applying (\ref{Doob}) iteratively for $j = 1,\dots,J$ we obtain that the probability that  the martingale $(M_t^{n,x})_{\frac{k-1}{n} \le t \le \frac{k}{n}}$ moves all the way from $x$ to $x + 2a$ is less than $2^{-J} < 2^{- n^{\frac{1}{4}}a}$. This gives the desired contradiction to our assumption that $(M_t^{n,x})_{\frac{k-1}{n} \le t \le \frac{k}{n}}$ does so with probability bigger than $\frac{\kappa}{n}$.






 
 


\section{Overview of related results in the literature} In the last sections we have focused on specific aspects of the theory of mimicking processes. 
In this final section we provide an overview  of related results in the literature and give some context for theorems discussed above. %For the convenience of the reader we will repeat some of the points already given at earlier stages. 


An early influential result is the work of Strassen \cite{St65} who established that there exists a submartingale with marginals $\mu_1, \ldots, \mu_N$ if and only these measures are increasing in the increasing convex order induced by increasing convex functions. He also proved that there exists a martingale with values in $\R^d$ and marginals $\mu_1, \ldots, \mu_N$  if and only if these measures are increasing in convex order. 
In \cite{Ke72, Ke73}, Kellerer managed to extend Strassen's result on the existence of submartingales to an arbitrary family of marginals. As an important particular case, this yields the existence of a mimicking martingale if the marginals increase in convex order.
Over time a number of authors have given new approaches to Kellerer's theorem see \cite{HiRo12, HiRoYo14, Lo08a, Lo08b, Lo08d, BeHuSt16, BeJu21}. As discussed extensively above the work Lowther adds substantial new developments. He  characterizes when the Markov martingale can be chosen to be continuous, as well as adding a clear-cut uniqueness part to  Kellerer's original result, complementing the formal uniqueness result of Dupire \cite{Du94}. 
We also recall from  above, that the question whether the natural extension of Kellerer's result to higher dimension holds true remains completely open. 

%A remarkably clear cut result on the existence and uniqueness of a continuous Markov Martingale was established by Lowther \cite{Lo08a, Lo08b, Lo08d}: Assume that $(\mu_t)_{t\in [0,1]}$ is increasing in convex order, weakly continuous in $t\in [0,1]$ and that each $\mu_t$ has convex support. Then there exists a unique strong Markov martingale $(X_t)_{t\in [0,1]}$ such that $X_t\sim \mu_t$ for $t\in [0,1]$. In fact the convex support assumption is not necessary if one replaces continuity of paths by an adequate assumption of weak continuity in probability. 

\medskip


Given a continuum of marginals which increase in convex order (and maybe satisfies additional technical conditions), different authors have provided specific constructions of (not necessarily Markovian) martingales that match these marginals. A main motivation stems from the calibration problem in mathematical finance. An additional goal has often been to give constructions that optimize particular functionals given the martingale and marginal  constraints since this yields robust bounds on option prices. Madan and Yor \cite{MaYo02} and 
K\''allblad, Tan, and Touzi \cite{KaTaTo15} establish a time continuous version of the Azema-Yor embedding. 
Hobson \cite{Ho16} established a continuous time version of the martingale  coupling constructed in \cite{HoKl13}.
Henry-Labordere, Tan, and Touzi \cite{HeTaTo16} as well as
Br\''uckerhoff, Huesmann, and Juillet \cite{BrHuJu20} provide continuous time versions of the shadow coupling (originally introduced in \cite{BeJu16}). Richard, Tan, and Touzi \cite{RiTaTo20} give a continuous time version of the Root solution to the Skorokhod embedding problem.
In a slightly different but related direction, Boubel and Juillet  \cite{BoJu18} consider a continuum  of marginals on the real line that do not satisfy an order condition and construct a canonical Markov-process matching these marginals.  We also refer to the book \cite{HiPrRoYo11} of Hirsch, Profeta, Roynette, and Yor that collects a variety of related constructions.  



The problem of finding martingales with given marginals has received specific attention in the case where these marginals equal the ones of Brownian motion. Hamza-Klebaner \cite{HaKl07}  posed the challenge of constructing martingales with Brownian marginals that differ from Brownian motion, so called \emph{fake Brownian motions}. 
 Non-continuous solutions can be found in  Madan-Yor \cite{MaYo02}, Hamza-Klebaner \cite{HaKl07}, Hobson \cite{Ho13}, and Fan-Hamza-Klebaner \cite{FaHaKl15} whereas continuous (but non-Markovian) fake Brownian motions were constructed by Oleszkiewicz \cite{Ol08}, Albin \cite{Al08}, Baker-Donati-Yor \cite{BaDoYo11} and Hobson \cite{Ho16}.
As already noted,  the accompanying article \cite{BePaSc21b} establishes that there exists  a Markovian martingale with continuous paths that  has Brownian marginals. 

\medskip 



A somewhat different direction arises if one starts with marginals that do not merely satisfy a structural condition (specifically,  monotonicity in  convex order) but rather assumes that a set of marginals is generated from an Ito-diffusion
\begin{align}\label{SourceSDE} dX_t = \sigma_t \, dB_t + \mu_t\, dt\end{align}
and one seeks a Markovian diffusion 
$$ d\hat X_t = \hat \sigma_t (X_t) \, dB_t + \mu_t(X_t)\, dt$$
that mimicks the evolution of $X$ in the sense that $X_t \sim \hat X_t$ for each $t\geq 0$. The process $\hat X$ is then considered to be the \emph{Markovian projection} of $X$. This line of research goes back essentially to the work of Krylov \cite{Kr85} and Gy\''ongy \cite{Gy86}. Of course, also the work of Dupire \cite{Du94} can be seen as a formal contribution to this line of research.
A rigorous justification of Dupire's formula under rather general assumption is obtained by Klebaner \cite{Kl02}. 
A very general theorem on mimicking aspects of Ito processes is given by Brunick and Shreve \cite{BrSh13}. Recently, Lacker, Shkolnikov, and Zhang \cite{LaShZh20} show that the results of \cite{Gy86, BrSh13} can be established directly from the superposition principle of Trevisan \cite{Tr16} (or Figalli \cite{Fi08} in the case where \eqref{SourceSDE} has bounded coefficients). (Notably, the main focus of the work  \cite{LaShZh20} is a mimicking result that shows that conditional time marginals of an Ito-process can be matched by a solution of a conditional McKean-Vlasov SDE with Markovian coefficients.) 


	
	%In contrast it is known from Gyöngy \cite{Gy86}, Dupire \cite{Du94}, and ultimately Lowther \cite{Lo08b} that Brownian motion is the unique continuous strong Markov martingale with Brownian marginals. 

In the mathematical finance community, Markovian (local vol) models are often considered to exhibit dynamics that are not particularly realistic. There has been significant interest to combine the convenience that the local vol model offers in terms of calibration with more realistic dynamics that are exhibited by other classes of financial models. That is,  given a Markovian model $d\hat X_t = \hat \sigma_t (\hat X_t)\, dS_t$ that represents marked data, one would like to ``reconstruct'' a more  realistic model $d X_t = \sigma_t \, dB_t$ and thus to ``invert'' the Markovian projection. A concrete way to perform this inversion is the \emph{stochastic local volatility model} see the work of Guyon and Henry-Labordere \cite{GuHe11, GuHe12, GuHe13}. However, it is remarkably delicate to establish existence and uniqueness results for the resulting SDEs.
Partial solutions where given by 
Jourdain and Zhou \cite{JoZh16} and by Lacker Shkolnikov and Zhang in  \cite{LaShZh19}. The problem is also discussed by Acciaio and Guyon \cite{AcGu20} who consider it an important open problem to establish existence of the stochastic local volatility model under fairly general assumptions. 



%\bibliographystyle{abbrv}
%\bibliography{joint_biblio}

\bibliographystyle{abbrv}
%\bibliography{joint_biblio, ../FilteredProcesses/lib/joint_biblio}
%\bibliography{../FilteredProcesses/lib/joint_biblio}
\bibliography{../MBjointbib/joint_biblio}

\end{document}	
